% Referee revision, 23 October 2026.
% Original file preserved as psh_exhaustion(1).tex.
% See the separate referee report and unified diff for the audit trail.
\documentclass[11pt]{amsart}

\usepackage[T1]{fontenc}
\usepackage{lmodern}
\usepackage{amssymb,mathtools}
\usepackage{microtype}
\usepackage{enumitem}
\usepackage{xcolor}
\usepackage[colorlinks=true,linkcolor=blue!55!black,citecolor=blue!55!black,urlcolor=blue!55!black]{hyperref}

\newtheorem{theorem}{Theorem}[section]
\newtheorem{proposition}[theorem]{Proposition}
\newtheorem{lemma}[theorem]{Lemma}
\newtheorem{corollary}[theorem]{Corollary}
\theoremstyle{definition}
\newtheorem{definition}[theorem]{Definition}
\theoremstyle{remark}
\newtheorem{remark}[theorem]{Remark}
\numberwithin{equation}{section}

\newcommand{\C}{\mathbb C}
\newcommand{\R}{\mathbb R}
\newcommand{\PP}{\mathbf P}
\newcommand{\OO}{\mathcal O}
\newcommand{\Tot}{\operatorname{Tot}}
\newcommand{\Bl}{\operatorname{Bl}}
\newcommand{\ddc}{dd^c}
\newcommand{\psh}{\mathrm{psh}}
\newcommand{\Levi}{\mathcal L}
\newcommand{\barp}{\bar\partial}
\newcommand{\supp}{\operatorname{supp}}
\newcommand{\base}{\operatorname{pr}}
\newcommand{\Uone}{\mathrm U(1)}

\title[Continuous and smooth plurisubharmonic exhaustions]{Continuous and Smooth
Plurisubharmonic Exhaustions:\\
Threefold Counterexamples and Partial Results for Surfaces}
\author{}
\date{}
\subjclass[2020]{Primary 32F17; Secondary 32Q15, 32L05}
\keywords{Plurisubharmonic exhaustion, K\"ahler manifold, semipositive line bundle,
ruled surface, Hironaka modification, Ueda theory, class VII surface, holonomy}

\begin{document}


\begin{abstract}
Starting from Koike's modified Grauert example, we construct a smooth
quasi-projective K\"ahler threefold with a locally Lipschitz
plurisubharmonic exhaustion but no $C^2$ plurisubharmonic exhaustion. A modification over one point gives a non-K\"ahler threefold
with the same properties. For divisor complements in compact complex
surfaces, we prove that continuous and smooth plurisubharmonic
exhaustions are equivalent to a boundary intersection-matrix condition
in two settings: every nonprojective compactifier, without a K\"ahler
assumption, and every smooth $\PP^1$-bundle with reduced simple normal
crossing boundary, without a degree restriction. Further affirmative
results treat disconnected projective boundaries and several log
surface cases. For smooth square-zero boundaries, character extension
and irregularity give sufficient conditions independent of continuous
weak completeness. Scoped holonomy arguments sharpen the remaining
neighborhood problem. The unrestricted implication in dimension two,
including the general K\"ahler and quasi-projective cases, is not proved.
\end{abstract}

\maketitle

\section{Introduction and main results}\label{sec:intro}

The regularity of a plurisubharmonic exhaustion matters even for
K\"ahler manifolds. This paper constructs a smooth quasi-projective
threefold admitting a continuous plurisubharmonic exhaustion but no
$C^2$ one, and then inserts a non-K\"ahler configuration over a single
point without changing these exhaustion properties. In dimension two,
the same line-bundle construction cannot give a counterexample. We
prove this obstruction and several affirmative regularization results
for surfaces, and specify the boundary problem left by our reductions.
Throughout the introduction, ``psh'' abbreviates plurisubharmonic, and
``smooth'' means $C^\infty$.

\subsection{Threefold counterexamples and their metric construction}

Our compact base is the modified Grauert ruled surface appearing in
Koike's work \cite[Example~4.1]{Koike2021}. Its failure of semipositivity
is already known; the first-order obstruction, including its $C^2$
version, is not claimed as a new result here
\cite[Theorem~1.4 and Example~4.1]{Koike2021},
\cite[Remark~2.8]{Wu2026}. We include a direct proof to make the
regularity threshold used below explicit. The analytic construction
then consists of a cusp correction, a global maximum, and a radius
criterion converting any hypothetical $C^2$ exhaustion into a
forbidden metric. Failure of smooth semipositivity by itself would
not provide a continuous exhaustion.

The cusp estimate also reaches the locally Lipschitz endpoint: a
sufficiently large multiple of $q^{1/2}$ dominates the initial
curvature in the scaled tangent variables. The coefficient is
essential at this endpoint; for smaller exponents shrinking the
neighborhood already suffices.

\begin{theorem}\label{thm:main}
Let $C$ be the smooth projective completion of
\[
 y^2=x^5-x,
\]
let $K=K_C$, and let
\begin{equation}\label{eq:extension-intro}
 0\longrightarrow K\longrightarrow V\longrightarrow\OO_C
 \longrightarrow0
\end{equation}
be a nonsplit holomorphic extension. Let
\[
 S=\PP_{\mathrm{lines}}(V),\qquad
 \pi:S\longrightarrow C,\qquad
 Y=\PP(K)\subset S,
\]
and set
\begin{equation}\label{eq:main-line}
 L=\OO_S(Y)\otimes\pi^*K,\qquad M=\Tot(L^{-1}).
\end{equation}
Then:
\begin{enumerate}[label=\textup{(\roman*)},leftmargin=2.5em]
 \item $S$ is a smooth projective surface and $M$ is a smooth
 quasi-projective K\"ahler threefold;
 \item $L$ admits a continuous Hermitian metric with plurisubharmonic
 local weights, but admits no $C^2$ Hermitian metric with semipositive
 curvature;
 \item $M$ admits a continuous plurisubharmonic exhaustion, but admits
 no $C^2$ plurisubharmonic exhaustion.
\end{enumerate}
The metric and the exhaustion in \textup{(ii)--(iii)} can be chosen
locally Lipschitz, and hence locally of class $C^{0,\gamma}$ for every
$0<\gamma\le1$. No assertion about existence of a $C^1$ exhaustion
is made.
\end{theorem}

The passage from metrics to exhaustions uses a criterion valid for a
holomorphic line bundle $A$ over any compact complex manifold $Z$:
\[
\begin{gathered}
 \Tot(A^{-1})\text{ admits a }C^r\text{ psh exhaustion}\\
 \Longleftrightarrow\\[-2pt]
 A\text{ admits a }C^r\text{ semipositive Hermitian metric},
 \qquad r\in\{2,\infty\}.
\end{gathered}
\]
This is Corollary~\ref{cor:criterion}. A semipositive metric gives the
squared dual norm exhaustion. Conversely, averaging an arbitrary
exhaustion under fiberwise circle rotation and taking a sufficiently
high level yields a radius function whose negative logarithm is a psh
metric weight. Consequently, the obstruction in Theorem~\ref{thm:main}
rules out every $C^2$ psh exhaustion, not only a smoothing of the
particular one constructed here.

The non-K\"ahler example has the same dimension and agrees with $M$
away from a compact exceptional fiber.

\begin{theorem}\label{thm:nonkahler-threefold}
For every point $x$ of the threefold $M$ in Theorem~\ref{thm:main}, there
is a smooth non-K\"ahler threefold $N$ and a proper holomorphic
modification $\mu:N\to M$ inducing a biholomorphism
\[
 N\setminus\mu^{-1}(x)\simeq M\setminus\{x\}.
\]
The threefold $N$ admits a locally Lipschitz psh exhaustion but no
$C^2$ psh exhaustion.
\end{theorem}

The construction places Hironaka's opposite-order blowups inside the
exceptional plane of the blowup of $x$. It produces two compact rational
curves whose sum is zero in real homology, precluding a K\"ahler form.
The exhaustion assertions follow from Proposition~\ref{prop:compact-modification}:
a proper surjective holomorphic map that is biholomorphic away from a compact
subset preserves the existence of $C^r$ psh exhaustions, for
$r=0,2,\infty$. The reverse implication uses convex composition to
make an exhaustion constant near the exceptional set before descending
it. Thus the modification is localized over one point; no small
deformation of a fixed complex structure is asserted. Choosing $x$
outside the zero section also preserves a neighborhood of the compact
base. The compact product principle in Proposition~\ref{prop:compact-product}
gives further examples in higher dimensions, including the Hopf-surface
product in Corollary~\ref{cor:nonkahler}.


\subsection{Surface results and the remaining boundary problem}

For a line bundle $A$ over a compact connected Riemann surface $B$,
the existence of a continuous psh exhaustion on $\Tot(A^{-1})$ is
equivalent both to $\deg A\ge0$ and to the existence of a smooth psh
exhaustion (Corollary~\ref{cor:curve-total-space}). Thus the threefold
construction cannot be lowered in dimension by replacing its compact
base with a curve.

Holomorphic convexity gives a different affirmative mechanism.
V\^aj\^aitu's theorem \cite{Vajaitu2022} implies that a connected
noncompact complex surface with a continuous psh exhaustion and a
nonconstant holomorphic function admits a smooth psh exhaustion
(Proposition~\ref{prop:surface-nonconstant-function}). A criterion of
Mongodi--Slodkowski also applies when an exhaustion takes arbitrarily
large values outside the image of the continuous minimal kernel
(Proposition~\ref{prop:surface-minimal-kernel}). Neither statement
requires a compactification. Their separate ambient real-analytic
regularization theorem has an additional hypothesis, recorded in
Subsection~\ref{subsec:surface-history}.

Our principal geometric results concern divisor complements
$X=S\setminus\supp D$, where $S$ is a smooth compact complex surface
and $D\ne0$ is effective. The existence of this compactification is an
additional assumption. Only the support of $D$ determines $X$.

\begin{theorem}[Two numerical criteria for surface exhaustions]
\label{thm:surface-criteria-intro}
Let $S$ be a connected smooth compact complex surface, let $D\ne0$
be effective, and put $X=S\setminus\supp D$. Suppose that either
\begin{enumerate}[label=\textup{(\roman*)},leftmargin=2.5em]
\item $S$ is nonprojective, with no K\"ahler or normal-crossing assumption; or
\item $S\to B$ is a smooth $\PP^1$-bundle over a compact connected
Riemann surface and $D$ is reduced with simple normal crossings.
\end{enumerate}
Then the following are equivalent:
\begin{enumerate}[label=\textup{(\alph*)},leftmargin=2.5em]
\item $X$ admits a continuous psh exhaustion;
\item $X$ admits a smooth psh exhaustion;
\item no connected component of $\supp D$ has negative definite
intersection matrix.
\end{enumerate}
In \textup{(ii)}, there is no restriction on the boundary degree or
the number of connected boundary components.
\end{theorem}

The two cases are proved in Theorem~\ref{thm:nonprojective} and
Corollary~\ref{cor:ruled-all-degrees}. The negative definite obstruction
is local: contraction to a normal surface point contradicts the
Hartogs maximum principle for a divergent psh function. For a
nonprojective surface of algebraic dimension one, the matrix condition
forces the boundary to be a union of whole elliptic fiber supports.
In algebraic dimension zero, reduction to a minimal model and Enoki's
theorem handle the class~VII case; the remaining Hopf boundaries are
smooth elliptic curves. The elliptic-neighborhood step uses the
September 2026 preprint \cite{CFSTZ2026}. In the smooth ruled case,
a connected square-zero boundary is either a fiber or a smooth
\'etale multisection. A finite \'etale base change reduces the latter
to an affine-line bundle, a punctured degree-zero line bundle, or a
product with a punctured projective line.

For projective compactifiers, disconnected boundary support also
implies the continuous-to-smooth conclusion
(Theorem~\ref{thm:surface-disconnected-boundary}). Resolution with the
full reduced boundary preserves the open surface. The intersection
matrix gives smooth strictly psh collars in the positive-eigenvalue
case. The residual projective case has a connected reduced simple
normal crossing boundary supporting a divisor
\[
 A=\sum_i a_iD_i,\quad a_i>0,\qquad
 A\text{ nef},\quad A^2=0,\quad A\cdot D_i=0,
\]
with $\kappa(S,A)=0$ and $K_S\cdot A\ge0$
(Corollary~\ref{cor:surface-connected-residual}). The continuous
exhaustion hypothesis is retained in this reduction.

We also prove the implication for smooth log Fano pairs, connected
numerically log Calabi--Yau pairs, and elliptic fibrations with vertical
boundary (Propositions~\ref{prop:surface-log-fano},
\ref{prop:surface-log-cy}, and \ref{prop:surface-elliptic-iitaka}).
These statements concern the indicated smooth compactifications;
arbitrary log minimal model steps need not preserve the complement.

For a smooth connected curve $C$ with $C^2=0$ in a compact K\"ahler
surface, Ueda theory supplies further sufficient conditions without
assuming continuous weak completeness
(Proposition~\ref{prop:surface-ueda-sufficient}). Let $\rho$ be the
unitary character of $N_{C/S}$ and put
\[
 K=\ker\bigl(H_1(C,\mathbb Z)\longrightarrow H_1(S,\mathbb Z)\bigr).
\]
If $\rho(K)$ is finite, a positive power of the character extends
to $S$, and $S\setminus C$ has a smooth psh exhaustion
(Theorem~\ref{thm:homology}). This holds in particular when the
homology map is injective, or when
$q(S)=h^1(S,\OO_S)\ge g(C)$
(Corollary~\ref{cor:irregularity}).

The examples in Section~\ref{sec:surface-examples} separate these
questions from holomorphic convexity and from the log Kodaira
dimension. Some complements have smooth psh exhaustions and only
constant holomorphic functions even when the ordinary canonical bundle
is ample. Other pairs realize arbitrary degree-zero normal bundles;
no continuous exhaustion is asserted for their arbitrary parameters.

The residual analytic problem already includes smooth connected
square-zero curves of genus at least two. Any counterexample of this
form must have $q(S)<g(C)$ and
\[
 \operatorname{rank}K=2\bigl(g(C)-q(S)\bigr),\qquad \rho(K)\text{ infinite},
\]
as well as an infinite-type neighborhood which is not unitary flat
(Corollary~\ref{cor:surface-residual-homology}). A continuous collar
which is constant on the leaves of a given holomorphic foliation
does force unitary linearization; a nonlinearizable irrational cyclic
holomorphic suspension admits no continuous divergent psh collar
(Propositions~\ref{prop:basic-collar} and
\ref{prop:cyclic-suspension}). These statements do not produce a
foliation or leafwise constancy from an arbitrary psh collar.
The full continuous-to-smooth implication for complex surfaces,
including all K\"ahler surfaces and all projective divisor complements,
remains unproved here.

Sections~\ref{sec:metrics}--\ref{sec:nonkahler} give the detailed
threefold constructions. Section~\ref{sec:surfaces} establishes the
surface reductions and the nonprojective criterion.
Section~\ref{sec:surface-log-cases} develops the homological, ruled,
and log surface results. Section~\ref{sec:surface-examples} gives the
examples, and Section~\ref{sec:surface-remaining} places the holonomy
obstructions beside the precise remaining problem.

\subsection{Relation to prior work and scope of the claims}
\label{subsec:provenance}

The known ruled-surface obstruction, Hironaka's opposite-order blowups,
and the quoted classification and Ueda theorems are inputs, not new
constructions of this paper. The nonconstant-holomorphic-function and
minimal-kernel criteria below are explicitly recorded consequences of
V\^aj\^aitu and Mongodi--Slodkowski. Likewise, the affine-line-bundle
step in the smooth-ruling argument reproduces a construction of
Koike--Ueda. The surface results combine these inputs with the boundary
and character arguments given here; we do not claim priority for every
intermediate criterion or for a general regularization theorem.

The self-contained threefold argument is the combination of the
explicit bounded-potential metric and the radius obstruction, followed
by a compactly supported modification. Remark~\ref{rem:minimal}
explains exactly which part of the bounded-potential question in
\cite[Section~4]{Zhang2025} this construction addresses. We make no
claim that dimension three is the smallest possible dimension of a
counterexample, nor that Lipschitz regularity is an optimal threshold.
The full dimension-two problem is not settled in this article.

Some surface conclusions use Theorem~1.6 and Proposition~6.2(B) of the
September 2026 preprint \cite{CFSTZ2026}; these uses are stated at the
relevant steps and are independent of the threefold construction.
The October 2026 preprint \cite{KoikeTakakura2026} concerns a different
neighborhood condition, as explained in
Remark~\ref{rem:pseudoconvex-neighborhoods}.

\subsection{Conventions and terminology}\label{sec:conventions}

We use
\[
 d^c=\frac{i}{4\pi}(\barp-\partial),\qquad
 \ddc=\frac{i}{2\pi}\partial\barp.
\]
For a holomorphic line bundle $A$ and a local holomorphic frame $e$, a
Hermitian metric is written
\[
 \|e\|_h^2=e^{-\varphi}.
\]
Its curvature is $\Theta_h(A)=\ddc\varphi$, representing $c_1(A)$.
If $e_j=g_{ij}e_i$, the metric weights satisfy
\begin{equation}\label{eq:metric-transition}
 \varphi_j=\varphi_i-\log|g_{ij}|^2.
\end{equation}
A positive continuous Hermitian metric is semipositive in the current
sense when its local weights are continuous psh functions. For a $C^2$
metric this agrees with pointwise semipositivity of its curvature.

For a $C^2$ real-valued function $f$, its unnormalized Levi form is
\[
 \Levi_f(\xi)=\sum_{j,k}f_{z_j\bar z_k}\xi_j\overline{\xi_k}.
\]
Thus $f$ is psh exactly when $\Levi_f\ge0$. The fixed positive factor
$1/(2\pi)$ in $\ddc$ does not affect any positivity argument.
Poincar\'e--Lelong takes the form
\begin{equation}\label{eq:PL}
 \ddc\log\|s_D\|_h^2=[D]-\Theta_h(\OO(D)),
\end{equation}
where $s_D$ is the canonical section of an effective divisor $D$.

\begin{definition}
A continuous real-valued function $F$ on a complex manifold $X$ is an
\emph{exhaustion} if every sublevel set $\{F\le c\}$ is compact.
Equivalently, $F(x_j)\to+\infty$ for every sequence escaping every
compact subset of $X$.
\end{definition}

We say that $X$ is \emph{continuously weakly complete} or
\emph{smoothly weakly complete} when it admits a continuous or smooth
psh exhaustion, respectively. Specifying the regularity avoids the
varying conventions for weak completeness in the literature.
Our projectivization $\PP_{\mathrm{lines}}(E)$ parametrizes
one-dimensional subspaces of the fibers of $E$, and its tautological
line bundle is $\OO_{\PP(E)}(-1)$.

\section{The ruled surface and continuous semipositive metrics}\label{sec:metrics}

\subsection{The ruled surface and the metric obstruction}\label{sec:ruled}

\subsubsection{The curve and the extension}

The polynomial $x^5-x$ has five distinct roots. The double covering
$C\to\PP^1$ is therefore branched at these five roots and at infinity.
Riemann--Hurwitz gives
\[
 2g(C)-2=2(-2)+6=2,
\]
so $g(C)=2$ and $\deg K=2$. Serre duality gives
\[
 \operatorname{Ext}^1(\OO_C,K)=H^1(C,K)
 \simeq H^0(C,\OO_C)^*\simeq\C.
\]
Consequently a nonzero extension class exists. All nonzero classes give
isomorphic extensions after allowing scalar automorphisms of the end
terms, so the isomorphism class of $V$ is unambiguous.

One can also specify the extension through a Dolbeault operator. Choose
a smooth positive $(1,1)$-form $\beta$ on $C$ with
$\int_C\beta=1$, and regard it as a $K$-valued $(0,1)$-form. On the
smooth bundle $K\oplus\OO_C$, set
\begin{equation}\label{eq:dolbeault-extension}
 \barp_V(s,f)=(\barp_Ks+f\beta,\barp f).
\end{equation}
This operator is integrable, since there are no $(0,2)$-forms on a
curve, and defines \eqref{eq:extension-intro}. A splitting would give a
holomorphic lift $(s,1)$ and hence $\barp_Ks=-\beta$. Since $s$ is a
smooth $(1,0)$-form, Stokes' theorem would then imply
$\int_C\beta=0$, a contradiction. The choice of the standard
identification of $K$-valued $(0,1)$-forms with $(1,1)$-forms only changes
this argument by a sign.

\subsubsection{The negative section and the line bundle}

By GAGA, $V$ is algebraic, so $S=\PP_{\mathrm{lines}}(V)$ is a smooth
projective surface. The inclusion $K\hookrightarrow V$ gives a section
$Y=\PP(K)$ of $\pi$. Its normal bundle is
\begin{equation}\label{eq:normal}
 N_{Y/S}\simeq\operatorname{Hom}(K,V/K)\simeq K^{-1}.
\end{equation}
In particular,
\begin{equation}\label{eq:restriction}
 \OO_S(Y)|_Y\simeq K^{-1},\qquad
 L|_Y\simeq\OO_Y,\qquad Y^2=-2.
\end{equation}
The trivialization of $L|_Y$ in \eqref{eq:restriction} is the canonical
one induced by $K^{-1}\otimes K$.

For context, this line bundle is nef and big. If $F$ is a fiber of $\pi$,
then
\[
 L\cdot Y=0,\qquad L\cdot F=1,\qquad L^2=2.
\]
For every irreducible curve $D\ne Y$ dominating $C$,
\[
 L\cdot D=Y\cdot D+2\deg(\pi|_D)>0.
\]
Every vertical irreducible curve is a fiber. Hence $L$ is nef, and the
usual surface criterion $L^2>0$ implies bigness. For completeness, the
latter also follows from Riemann--Roch: $h^2(S,L^m)=0$ for $m$ large,
because $(K_S-mL)$ has negative intersection with a fixed ample divisor,
and therefore
\[
 h^0(S,L^m)\ge\chi(S,L^m)=m^2+O(m).
\]
Moreover, $2L-Y=Y+2\pi^*K$ is ample by Nakai--Moishezon: its square is
$6$, its intersections with $Y$ and a fiber are $2$ and $1$, and its
intersection with every other irreducible curve is positive.

\subsubsection{First-order transition functions}

Choose a finite coordinate cover $\{U_i\}$ of $C$, frames $k_i$ of $K$,
and holomorphic lifts $v_i$ of $1$ in \eqref{eq:extension-intro}.
On overlaps,
\begin{equation}\label{eq:vector-transitions}
 k_j=g_{ij}k_i,\qquad v_j=v_i+c_{ij}k_i.
\end{equation}
The cocycle $\{c_{ij}k_i\}$ represents the nonzero extension class.
Near $Y$, write a line in $V$ as
\[
 [k_i+w_iv_i].
\]
Expressing this line in the $j$-th frame yields
\begin{equation}\label{eq:normal-transition}
 w_j=\frac{g_{ij}w_i}{1-c_{ij}w_i}.
\end{equation}
Shrinking the neighborhoods if necessary, all denominators are nonzero.
Let $d_i$ be the frame of $\OO_S(Y)$ defined by $s_Y=w_id_i$. Then
\[
 \frac{d_j}{d_i}=\frac{w_i}{w_j}
 =\frac{1-c_{ij}w_i}{g_{ij}}.
\]
Thus the frames $\ell_i=d_i\otimes\pi^*k_i$ of $L$ satisfy
\begin{equation}\label{eq:L-transition}
 \ell_j=(1-c_{ij}w_i)\ell_i.
\end{equation}
Their restrictions to $Y$ agree, in accordance with the trivialization
in \eqref{eq:restriction}.

\begin{proposition}\label{prop:no-metric}
The line bundle $L$ admits no $C^2$ Hermitian metric with semipositive
curvature.
\end{proposition}

\begin{proof}
Suppose such a metric exists, with $C^2$ psh weights $\varphi_i$ in the
frames $\ell_i$. By \eqref{eq:L-transition}, the functions
$\varphi_i|_Y$ agree on overlaps. They define a subharmonic function on
the compact connected curve $Y$, hence a constant. Rescaling the metric,
we may assume
\begin{equation}\label{eq:flat-value}
 \varphi_i(z,0)=0.
\end{equation}
In particular $\varphi_{i,z\bar z}(z,0)=0$. The Levi matrix at $(z,0)$
is positive semidefinite. A positive semidefinite Hermitian matrix with
a zero diagonal entry has zero entries in the corresponding row and
column. Therefore
\[
 \varphi_{i,w\bar z}(z,0)=0.
\]
The $C^1$ function
\[
 a_i(z)=\varphi_{i,w}(z,0)
\]
satisfies the Cauchy--Riemann equation and is holomorphic.

Equations \eqref{eq:metric-transition} and \eqref{eq:L-transition} give
\[
 \varphi_j(z,w_j)
 =\varphi_i(z,w_i)-\log|1-c_{ij}(z)w_i|^2,
\]
where the base coordinate change is understood. Differentiating with
respect to $w_i$ at zero and using \eqref{eq:normal-transition}, we obtain
\begin{equation}\label{eq:coboundary}
 g_{ij}a_j=a_i+c_{ij}.
\end{equation}
The holomorphic lifts $v_i'=v_i-a_ik_i$ consequently agree on overlaps:
\[
 v_j'=v_i+c_{ij}k_i-a_jg_{ij}k_i=v_i-a_ik_i=v_i'.
\]
They define a global holomorphic lift of $1$, splitting
\eqref{eq:extension-intro}. This is a contradiction.
\end{proof}

This calculation is the specialized first-order obstruction occurring
in \cite[Theorem~1.4 and Example~4.1]{Koike2021}. Related obstruction
classes, including their applicability to $C^2$ metrics, are discussed
in \cite{Wu2026}. Neither reference is needed for the preceding proof.

\subsection{A local continuous metric near a negative curve}\label{sec:cusp}

We isolate the analytic estimate used in the construction. The unitary
flatness assumption means that a smooth metric of zero curvature exists
on the restricted line bundle.

\begin{lemma}[Cusp correction]\label{lem:cusp}
Let $D$ be a compact connected smooth curve in a complex surface $X$. Suppose
$\deg N_{D/X}<0$, and let $A$ be a holomorphic line bundle such that
$A|_D$ is unitary flat. Choose a smooth metric $h_0$ on $A$ restricting
to a flat metric on $D$, and a smooth metric $h_D$ on $\OO_X(D)$ whose
curvature is strictly negative on $TD$. Put
\[
 q=\|s_D\|_{h_D}^2.
\]
For every $0<a\le\tfrac12$, there is a constant $\Lambda>0$ such that
the local weights of $h_0$ plus $\Lambda q^a$ are strictly psh off $D$
on a sufficiently small neighborhood of $D$ and extend as continuous
psh weights across $D$. They are locally of class $C^{0,2a}$; in
particular they are locally Lipschitz when $a=\tfrac12$.
For $0<a<\tfrac12$ one may take $\Lambda=1$ after shrinking the
neighborhood. At $a=\tfrac12$, a sufficiently large $\Lambda$ suffices.
\end{lemma}

\begin{proof}
The chosen metrics exist. Negative degree on the curve gives a negative
curvature metric on $N_{D/X}=\OO_X(D)|_D$. Any metric on a restricted
line bundle extends smoothly to a neighborhood: extend the logarithm
of its ratio to a fixed smooth metric. The same argument extends the
flat metric on $A|_D$.

Use finitely many relatively compact coordinates $(z,w)$ near $D$,
with $D=\{w=0\}$. Write
\[
 q=|w|^2e^{-b(z,w)},\qquad \beta=-b,
\]
where $b$ is the smooth weight of $h_D$. Negativity along $TD$ says
\begin{equation}\label{eq:beta-positive}
 \beta_{z\bar z}(z,0)>0.
\end{equation}
For $w\ne0$, set $\tau=\log q=\log|w|^2+\beta$. The chain rule gives
the exact identity
\begin{equation}\label{eq:cusp-exact}
 \Levi_{q^a}(v)
 =q^a\bigl(a\Levi_{\tau}(v)+a^2|\partial\tau(v)|^2\bigr).
\end{equation}
Evaluate this on
\[
 v=\xi\frac{\partial}{\partial z}
    +w\nu\frac{\partial}{\partial w}.
\]
Since $\log|w|^2$ is pluriharmonic off $D$,
\[
 \Levi_{\tau}(v)
 =\beta_{z\bar z}(z,0)|\xi|^2
   +O(|w|)(|\xi|^2+|\nu|^2),
\]
and
\[
 \partial\tau(v)
 =\nu+\beta_z(z,0)\xi+O(|w|)(|\xi|+|\nu|).
\]
Therefore
\begin{align}\label{eq:scaled-cusp}
 \frac{\Levi_{q^a}(\xi,w\nu)}{q^a}
 &=a\beta_{z\bar z}(z,0)|\xi|^2
   +a^2|\nu+\beta_z(z,0)\xi|^2\notag\\
 &\hspace{1em}+O(|w|)(|\xi|^2+|\nu|^2).
\end{align}
All remainder bounds are uniform after choosing the finite cover and
shrinking its coordinate neighborhoods.

The leading Hermitian form in \eqref{eq:scaled-cusp} is positive
definite: vanishing of its first term forces $\xi=0$ by
\eqref{eq:beta-positive}, and vanishing of the second then forces
$\nu=0$. Its smallest eigenvalue has a positive uniform lower bound on
the compact coordinate pieces covering $D$. Also $q^a$ is uniformly
comparable to $|w|^{2a}$. Thus, after shrinking the neighborhoods,
\begin{equation}\label{eq:cusp-lower}
 \Levi_{q^a}(\xi,w\nu)
 \ge c|w|^{2a}(|\xi|^2+|\nu|^2)
\end{equation}
for a uniform $c>0$.

Let $\varphi_0$ be a local weight of $h_0$. Its restriction to $D$ has
zero curvature, so $\varphi_{0,z\bar z}(z,0)=0$. Smoothness gives
\[
 \varphi_{0,z\bar z}=O(|w|),\qquad
 \varphi_{0,z\bar w}=O(1),\qquad
 \varphi_{0,w\bar w}=O(1).
\]
Consequently, in the same scaled variables,
\begin{equation}\label{eq:baseline-bound}
 |\Levi_{\varphi_0}(\xi,w\nu)|
 \le C|w|(|\xi|^2+|\nu|^2).
\end{equation}
Combining \eqref{eq:cusp-lower} and \eqref{eq:baseline-bound} yields
\begin{equation}\label{eq:cusp-coefficient}
 \Levi_{\varphi_0+\Lambda q^a}(\xi,w\nu)
 \ge\bigl(\Lambda c|w|^{2a}-C|w|\bigr)
       (|\xi|^2+|\nu|^2).
\end{equation}
If $0<a<\tfrac12$, take $\Lambda=1$ and then shrink until the
coefficient is positive. If $a=\tfrac12$, fix $\Lambda>C/c$; the
coefficient becomes $(\Lambda c-C)|w|>0$ throughout the punctured
neighborhood where the two uniform estimates hold. The constants
$c,C$ were obtained before choosing $\Lambda$, so this choice is not
circular. Every tangent vector at $w\ne0$ can be written in these
scaled variables. This proves strict plurisubharmonicity off $D$ in
both cases.

The function $q^a=|w|^{2a}e^{-ab}$ is continuous and locally
$C^{0,2a}$, including $2a=1$. The local weights
$\varphi_0+\Lambda q^a$ are bounded near $D$ and extend continuously
to it. The removable singularity theorem for psh functions across an
analytic hypersurface therefore proves that these extended weights
are psh. Adding the globally defined function $\Lambda q^a$ preserves
the metric transition law \eqref{eq:metric-transition}.
\end{proof}

\begin{remark}
The scale in \eqref{eq:scaled-cusp} is essential. The normal component
of the Levi form has order $|w|^{2a-2}$, while the tangential component
has order $|w|^{2a}$. These two scales allow the cusp to dominate the
bounded mixed curvature of a smooth initial metric. At $a=\tfrac12$
the scaled cusp and the initial error both have order $|w|$, which
explains the coefficient choice in \eqref{eq:cusp-coefficient}.
Ordinary smoothing of the cusp would not preserve this domination at
$D$. The argument proves no $C^1$ regularity.
\end{remark}

\subsection{An explicit global continuous metric}\label{sec:global}

Return to the ruled surface $S$ and the line bundle $L$ in
\eqref{eq:main-line}. Choose a positive smooth metric $h_K$ on $K$.
Using \eqref{eq:restriction}, extend the dual metric $h_K^{-1}$ on
$\OO_S(Y)|_Y$ to a smooth metric $h_Y$ on $\OO_S(Y)$. Set
\begin{equation}\label{eq:baseline-metric}
 h_0=h_Y\otimes\pi^*h_K,\qquad
 \theta_0=\Theta_{h_0}(L),\qquad
 q=\|s_Y\|_{h_Y}^2.
\end{equation}
The metric $h_0|_Y$ is the flat metric under the canonical
trivialization of $L|_Y$, and
\begin{equation}\label{eq:curvature-restrictions}
 \theta_0|_{TY}=0,\qquad
 \Theta_{h_Y}(\OO(Y))|_{TY}=-\Theta_{h_K}(K)<0.
\end{equation}

Fix $0<a\le\tfrac12$ and choose $\Lambda>0$ as in
Lemma~\ref{lem:cusp}. Since $q$ vanishes exactly on $Y$ and $S$ is
compact, there exists $\varepsilon>0$ such that the weights
$\varphi_0+\Lambda q^a$ are psh on $\{q<\varepsilon\}$.
Write $Q=\max_Sq$ and choose $0<\varepsilon<Q$. Choose a constant
\begin{equation}\label{eq:B-choice}
 B>\max_{\varepsilon\le t\le Q}\bigl(\Lambda t^a-\log t\bigr).
\end{equation}
Define, on $[0,Q]$,
\begin{equation}\label{eq:chi}
 \chi(0)=0,\qquad
 \chi(t)=\max\{\Lambda t^a,\log t+B\}\quad(t>0),
\end{equation}
and put
\begin{equation}\label{eq:global-metric}
 h=h_0\exp\bigl(-\chi(q)\bigr).
\end{equation}

\begin{proposition}\label{prop:continuous-metric}
The metric \eqref{eq:global-metric} is a positive continuous Hermitian
metric on $L$ with psh local weights. It is locally $C^{0,2a}$,
and is locally Lipschitz for $a=\tfrac12$.
\end{proposition}

\begin{proof}
The logarithmic branch in \eqref{eq:chi} tends to $-\infty$ as $t\to0$,
so $\chi(t)=\Lambda t^a$ for all sufficiently small $t$. Hence $\chi$ is
continuous at zero and bounded on $[0,Q]$.

On $S\setminus Y$, Poincar\'e--Lelong and
\eqref{eq:baseline-metric} give
\begin{equation}\label{eq:log-curvature}
 \theta_0+\ddc\log q=\pi^*\Theta_{h_K}(K)\ge0.
\end{equation}
As a current on $S$, the same expression equals
$\pi^*\Theta_{h_K}(K)+[Y]$. Thus $\varphi_0+\log q+B$ is a singular
psh weight of $L$.

On $\{0<q<\varepsilon\}$, the weights of $h$ are
\[
 \varphi_0+\chi(q)
 =\max\{\varphi_0+\Lambda q^a,\varphi_0+\log q+B\},
\]
a maximum of two psh functions. They are therefore psh. On
$\{q\ge\varepsilon\}$, \eqref{eq:B-choice} makes the logarithmic
branch strictly dominant. This domination also holds on a collar of
$\{q=\varepsilon\}$, so no boundary gluing issue arises. The weights
are psh there by \eqref{eq:log-curvature}. Near $Y$, the cusp branch is
dominant, and Lemma~\ref{lem:cusp} proves positivity across $Y$.

All the weights are finite and continuous. Near $Y$ their only
nonsmooth term is $\Lambda q^a=\Lambda|w|^{2a}e^{-ab}$,
which is $C^{0,2a}$ for $0<2a\le1$.
Away from $Y$, both branches are smooth, and their maximum is locally
Lipschitz, hence locally $C^{0,2a}$. This proves the asserted regularity.
\end{proof}

\begin{remark}[Bounded potentials and minimal singularities]\label{rem:minimal}
The current
\[
 T=\theta_0+\ddc\chi(q)\in c_1(L)
\]
has locally Lipschitz potentials when $a=\tfrac12$. It also has minimal
singularities: if $T'=\theta_0+\ddc v\ge0$ is another positive current
in $c_1(L)$, its quasi-psh potential $v$ is bounded above on compact
$S$, while $\chi(q)$ is bounded below. Hence
\[
 v\le\chi(q)+C_v
\]
for a constant $C_v$. This is precisely the comparison expressing that
$T$ is no more singular than $T'$. Here global potentials exist by the
$\partial\bar\partial$-lemma on the projective surface.

For this modified Grauert example, Proposition~\ref{prop:continuous-metric}
therefore gives an affirmative answer to the bounded-potential question
in \cite[Section~4]{Zhang2025} for the modified example. That question
also mentions Grauert's original example; the present argument does
not identify its line bundle with $L$ and does not settle that other
case. The metric is constructed here, not taken from that preprint.
This explicit answer for $L$ is not a statement about arbitrary nef
and big line bundles, nor a claim that every minimal-singularity
representative has continuous potentials.
\end{remark}

\section{Line-bundle exhaustions and the K\"ahler counterexample}\label{sec:radius}

\subsection{Dual norms and compact-group averaging}

\begin{lemma}\label{lem:norm}
Let $Z$ be a compact complex manifold and $A\to Z$ a holomorphic line
bundle with a positive continuous Hermitian metric $h$ having psh
local weights. Then the squared dual norm
\[
 \rho(v)=\|v\|_{h^{-1}}^2\qquad(v\in\Tot(A^{-1}))
\]
is a continuous psh exhaustion. If the weights are locally
$C^{0,\gamma}$ for $0<\gamma\le1$, so is $\rho$.
\end{lemma}

\begin{proof}
For a local frame $e_i$ of $A$, write a dual vector as
$v=\lambda_i e_i^*$, and write $h(e_i,e_i)=e^{-\varphi_i}$. Then
\begin{equation}\label{eq:local-norm}
 \rho=|\lambda_i|^2e^{\varphi_i}.
\end{equation}
Off the zero section, $\log\rho=\log|\lambda_i|^2+\varphi_i$ is psh.
Composition with the increasing convex function $\exp$ makes $\rho$
psh there. Formula \eqref{eq:local-norm} extends continuously by zero
to the zero section and is locally bounded, so removability gives psh
on the whole total space.

Fix a smooth Hermitian metric $h_{\mathrm{ref}}$ on $A^{-1}$. By
compactness of $Z$ and positivity of $h$, constants $c,C>0$ satisfy
\[
 c\|v\|_{h_{\mathrm{ref}}}^2
 \le\rho(v)\le C\|v\|_{h_{\mathrm{ref}}}^2.
\]
The closed disk bundle of any fixed radius over $Z$ is compact.
Every sublevel of $\rho$ is a closed subset of such a disk bundle and
is therefore compact. The H\"older assertion follows locally from
\eqref{eq:local-norm}.
\end{proof}

\begin{lemma}[Averaging preserves exhaustions]\label{lem:average}
Let a compact group $G$ act continuously on a locally compact Hausdorff
space $X$, and let $F$ be a continuous exhaustion. Its normalized Haar average
\[
 \overline F(x)=\int_G F(gx)\,dg
\]
is also an exhaustion. If $G$ is a compact Lie group, $X$ is a complex
manifold, the action is smooth and holomorphic, and $F$ is $C^2$ psh,
then $\overline F$ is $C^2$ psh. The same assertion holds in the
smooth category.
\end{lemma}

\begin{proof}
For $K_c=\{F\le c\}$, the orbit saturation $G\cdot K_c$ is compact.
If $x\notin G\cdot K_c$, then $F(gx)>c$ for every $g\in G$, and
continuity on the compact orbit implies $\overline F(x)>c$. Hence
\[
 \{\overline F\le c\}\subset G\cdot K_c.
\]
The sublevel on the left is closed and therefore compact. Differentiation
under the integral preserves $C^2$ regularity, and averaging the
nonnegative Levi forms of the holomorphic pullbacks preserves psh.
\end{proof}

\subsection{The radius criterion}

\begin{proposition}[The radius criterion]\label{prop:radius}
Let $Z$ be a compact complex manifold and $A\to Z$ a holomorphic line
bundle. If $\Tot(A^{-1})$ admits a $C^2$ psh exhaustion, then $A$ admits
a $C^2$ Hermitian metric with semipositive curvature. If the exhaustion
is smooth, the metric is smooth.
\end{proposition}

\begin{proof}
Average the exhaustion over the holomorphic circle action
$v\mapsto e^{i\alpha}v$ on the fibers, using
Lemma~\ref{lem:average}. Denote the resulting invariant exhaustion by
$F$. For a local frame $e_i$ of $A$ and $v=\lambda_i e_i^*$, circle
invariance gives
\begin{equation}\label{eq:radial-function}
 F(z,\lambda_i)=U_i(z,t),\qquad t=\log|\lambda_i|^2,
 \quad\lambda_i\ne0.
\end{equation}
The function
\begin{equation}\label{eq:holomorphic-pullback}
 V_i(z,\zeta)=U_i(z,\operatorname{Re}\zeta)
             =F(z,e^{\zeta/2})
\end{equation}
is a holomorphic pullback of $F$ and is psh. In particular
$U_{i,tt}\ge0$.

For fixed $z$, continuity at the zero section gives the finite limit
\[
 \lim_{t\to-\infty}U_i(z,t)=F(z,0).
\]
Exhaustiveness gives $U_i(z,t)\to+\infty$ as $t\to+\infty$.
A differentiable convex function with finite left limit has
nonnegative derivative. Moreover, if its derivative vanishes at
$t_0$, convexity forces it to be constant on $(-\infty,t_0]$.
Consequently
\begin{equation}\label{eq:radial-positive}
 U_{i,t}(z,t)>0\quad\text{whenever }U_i(z,t)>F(z,0).
\end{equation}

Fix $c>\max_ZF|_Z$. By convexity, the two limits, and
\eqref{eq:radial-positive}, there is a unique solution
\begin{equation}\label{eq:implicit-radius}
 U_i(z,\sigma_i(z))=c,
 \qquad U_{i,t}(z,\sigma_i(z))>0.
\end{equation}
The implicit function theorem makes $\sigma_i$ of class $C^2$, or
smooth when $F$ is smooth.

We verify the sign of its Levi form explicitly. Differentiating
\eqref{eq:implicit-radius} twice in base coordinates gives
\begin{align}\label{eq:implicit-hessian}
 0={}&U_{i,z_j\bar z_k}
 +U_{i,z_jt}\sigma_{i,\bar z_k}
 +U_{i,t\bar z_k}\sigma_{i,z_j}\notag\\
 &+U_{i,tt}\sigma_{i,z_j}\sigma_{i,\bar z_k}
 +U_{i,t}\sigma_{i,z_j\bar z_k},
\end{align}
with all derivatives of $U_i$ evaluated at $(z,\sigma_i(z))$.
For $V_i$ in \eqref{eq:holomorphic-pullback}, its $\zeta\bar\zeta$
entry is $U_{i,tt}/4$, and its mixed entries have a factor $1/2$.
Thus \eqref{eq:implicit-hessian} is precisely the identity
\begin{equation}\label{eq:radius-Levi}
 \Levi_{V_i}\bigl(\xi,2\partial\sigma_i(\xi)\bigr)
 =-U_{i,t}\Levi_{\sigma_i}(\xi).
\end{equation}
The left side is nonnegative and $U_{i,t}>0$, so
$\psi_i=-\sigma_i$ is psh.

If $e_j=g_{ij}e_i$, the dual coordinates satisfy
$\lambda_j=g_{ij}\lambda_i$. The level in
\eqref{eq:implicit-radius} consequently transforms by
\[
 \sigma_j=\sigma_i+\log|g_{ij}|^2,
\]
so
\[
 \psi_j=\psi_i-\log|g_{ij}|^2.
\]
These are exactly the metric transition laws
\eqref{eq:metric-transition}. Hence $\|e_i\|^2=e^{-\psi_i}$ defines
the desired semipositive Hermitian metric on $A$.
\end{proof}

\begin{corollary}\label{cor:criterion}
For a holomorphic line bundle $A$ over a compact complex manifold,
$\Tot(A^{-1})$ admits a $C^2$ psh exhaustion if and only if $A$ admits
a $C^2$ semipositive Hermitian metric. The same equivalence holds with
$C^2$ replaced by $C^\infty$.
\end{corollary}

\begin{proof}
The forward implication is Proposition~\ref{prop:radius}. For the
reverse implication, the squared dual norm in Lemma~\ref{lem:norm}
has the regularity of the given metric, by its local expression
\eqref{eq:local-norm}.
\end{proof}

\subsection{Proof of the K\"ahler counterexample}\label{sec:conclusion}

\begin{proof}[Proof of Theorem~\ref{thm:main}]
The construction in Subsection~\ref{sec:ruled} gives a smooth projective
surface $S$ and an algebraic line bundle $L$. The total space $M$ is
the open subset
\begin{equation}\label{eq:compactification}
 M=\PP_{\mathrm{lines}}(\OO_S\oplus L^{-1})
       \setminus\PP_{\mathrm{lines}}(L^{-1}).
\end{equation}
The identification sends $v\in L^{-1}_z$ to the line $[(1,v)]$.
The projective bundle in \eqref{eq:compactification} is a smooth
projective threefold. Restricting any of its K\"ahler forms proves
that $M$ is K\"ahler.

Let $h$ be the continuous semipositive metric from
Proposition~\ref{prop:continuous-metric}. If $p:M\to S$ is the bundle
projection, the exhaustion is explicitly
\begin{equation}\label{eq:explicit-exhaustion}
 \rho(v)=\|v\|_{h^{-1}}^2
       =\|v\|_{h_0^{-1}}^2
          \exp\!\bigl(\chi(q(p(v)))\bigr).
\end{equation}
Lemma~\ref{lem:norm} shows that $\rho$ is continuous, psh, and
exhaustive. Taking $a=\tfrac12$ and the coefficient $\Lambda$ supplied
by Lemma~\ref{lem:cusp} makes $h$, its weights, and $\rho$ locally
Lipschitz. In particular the same example has local $C^{0,\gamma}$
regularity for all $0<\gamma\le1$.

If a $C^2$ psh exhaustion on $M$ existed,
Proposition~\ref{prop:radius} would produce a $C^2$ semipositive
Hermitian metric on $L$. Proposition~\ref{prop:no-metric} excludes
such a metric. Therefore no $C^2$ psh exhaustion exists, completing
the proof.
\end{proof}

\begin{remark}[Interpretation]
The example separates existence of a continuous exhaustion from
existence of any smooth exhaustion. The conclusion does not depend on
whether the particular function \eqref{eq:explicit-exhaustion} can be
smoothed: every hypothetical $C^2$ exhaustion leads to the forbidden
metric. The argument does not determine whether a $C^1$ psh
exhaustion exists, and the surface results below do not establish that
dimension three is minimal. Thus a definition of weak pseudoconvexity
requiring a smooth
psh exhaustion is strictly stronger than the corresponding definition
requiring only a continuous psh exhaustion, even for smooth
quasi-projective K\"ahler threefolds.
\end{remark}

\section{Compact modifications and non-K\"ahler counterexamples}
\label{sec:nonkahler}

The K\"ahler counterexample can be changed over a single point to obtain
a non-K\"ahler counterexample of the same dimension.  We first isolate
the exhaustion argument, then insert Hironaka's opposite-order blowup
configuration in an exceptional plane.  Compact products give a simpler
alternative when an increase in dimension is allowed.

\subsection{Exhaustions under compactly supported modifications}

\begin{proposition}\label{prop:compact-modification}
Let $f:X'\to X$ be a proper surjective holomorphic map between nonempty
complex manifolds.  Suppose that, for a compact set $K\subset X$, it
restricts to a biholomorphism
\[
 X'\setminus f^{-1}(K)\longrightarrow X\setminus K.
\]
For every $r\in\{0,2,\infty\}$, $X$ admits a $C^r$ psh exhaustion
if and only if $X'$ does.
\end{proposition}

\begin{proof}
The case $K=\varnothing$ is immediate.  Holomorphic pullback preserves
plurisubharmonicity and the stated
regularity, and properness makes the sublevels of a pulled-back
exhaustion compact.  Conversely, suppose $F$ is a $C^r$ psh exhaustion
on $X'$.  Choose a relatively compact open neighborhood $U$ of $K$ and
\[
 c>\max_{f^{-1}(\overline U)}F.
\]
Choose a smooth convex nondecreasing function
$\chi:\R\to[0,\infty)$ that vanishes on $(-\infty,c]$ and tends to
infinity.  For example, integrate a smooth nondecreasing function
which is zero up to $c$ and one from $c+1$ onward.  On $X\setminus K$
define
\[
 G=\chi\circ F\circ(f|_{X'\setminus f^{-1}(K)})^{-1}.
\]
This function is $C^r$ and psh.  It vanishes on $U\setminus K$, so
extension by zero on $K$ gives a $C^r$ psh function on all of $X$.
For each nonempty sublevel $\{G\le a\}$ choose $b$ such that
$\chi(t)\le a$ implies $t\le b$.  Then
\[
 \{G\le a\}\subset K\cup f(\{F\le b\}).
\]
The left side is closed and the right side is compact.  Thus $G$ is an
exhaustion; no smooth descent of the original function $F$ is needed.
\end{proof}

\subsection{A non-K\"ahler modification over one point}

\begin{proposition}\label{prop:point-modification}
Let $X$ be a connected smooth complex threefold and $x\in X$.
There is a connected smooth non-K\"ahler threefold $N$ and a proper
holomorphic modification $\mu:N\to X$ inducing a biholomorphism
\[
 N\setminus\mu^{-1}(x)\simeq X\setminus\{x\}.
\]
\end{proposition}

\begin{proof}
Blow up $x$, obtaining
\[
 b:W=\Bl_xX\longrightarrow X,\qquad E=b^{-1}(x)\simeq\PP^2.
\]
In homogeneous coordinates $[u:v:w]$ on $E$, take
\[
 A=\{w=0\},\qquad B=\{uv=w^2\}.
\]
The line $A$ and smooth conic $B$ meet transversely at exactly
$p=[1:0:0]$ and $q=[0:1:0]$.  Near either point they are the two
coordinate axes in a smooth surface $E$ inside $W$.

Over $W\setminus\{q\}$, first blow up $A\setminus\{q\}$ and then
the strict transform of $B\setminus\{q\}$.  Over
$W\setminus\{p\}$, first blow up $B\setminus\{p\}$ and then the
strict transform of $A\setminus\{p\}$.  All these centers are closed
smooth curves in the respective open manifolds.  On
$W\setminus\{p,q\}$ the centers are disjoint, so the blowups commute
canonically over the base.  Glue the two resulting manifolds through
this isomorphism and write
\[
 \nu:N\longrightarrow W,\qquad \mu=b\circ\nu.
\]
The glued space is Hausdorff: points with distinct images can be
separated in $W$, and points with the same image belong to one of the
Hausdorff open manifolds lying over the two base opens.  It is smooth
of dimension three because every blowup center is smooth.  Properness
is local on the target, so $\nu$ and $\mu$ are proper.  Each blowup
chart is connected and the charts overlap, hence $N$ is connected.
All centers lie in $E$, so $\mu$ is an isomorphism away from $x$.
This is Hironaka's opposite-order construction
\cite[Section~6]{Barlet2015}; see also \cite[Tag~08J1]{StacksHironaka}.

For completeness, the obstruction to a K\"ahler form is entirely
contained in compact exceptional surfaces.  Let $D_A,D_B\subset N$
be the exceptional divisors dominating $A,B$, and let
$\alpha,\beta\in H_2(N,\R)$ be the classes of their general
$\PP^1$ fibers.  The surfaces $D_A,D_B$ are smooth and compact:
locally each is a ruled surface or its blowup at a point, and they
map properly to the compact curves $A,B$.  The fiber divisors of
each map are linearly equivalent, since $A\simeq B\simeq\PP^1$,
and hence are homologous on the corresponding compact surface.

Over $p$, the strict transform of $B$ meets the first exceptional
divisor transversely at one point.  The second blowup therefore
restricts on that divisor to a surface blowup at that point.  Its
fiber over $p$ is the reduced union of the strict transform
$R_p\simeq\PP^1$ of the first exceptional fiber and the new
exceptional fiber $T_p\simeq\PP^1$. Both occur with multiplicity one.
Indeed, take local coordinates $(x,y,z)$ with
$A=\{y=z=0\}$ and $B=\{x=z=0\}$. In the first blowup chart
$z=yt$, the first exceptional divisor is $y=0$ and the second center
is $x=t=0$. Its restriction to $y=0$ is the blowup of the smooth
surface with coordinates $(x,t)$ at the origin. The pullback of the
fiber equation $x=0$ is its strict transform plus the exceptional
curve, each once. The curve $T_p$ is also the fiber of $D_B\to B$
over $p$. Thus
\[
 \alpha=[R_p]+[T_p],\qquad \beta=[T_p].
\]
Over $q$ the order is reversed.  With analogous notation,
\[
 \beta=[R_q]+[T_q],\qquad \alpha=[T_q].
\]
These relations hold in ordinary homology: they arise from fiber
divisors on compact $D_A$ and $D_B$, regardless of compactness of $N$.
Adding them gives
\begin{equation}\label{eq:hironaka-zero-cycle}
 [R_p]+[R_q]=0\quad\text{in }H_2(N,\R).
\end{equation}
If $\omega$ were a K\"ahler form on $N$, closedness and
\eqref{eq:hironaka-zero-cycle} would imply
\[
 0=\int_{R_p}\omega+\int_{R_q}\omega>0,
\]
a contradiction.
\end{proof}

\begin{proof}[Proof of Theorem~\ref{thm:nonkahler-threefold}]
Apply Proposition~\ref{prop:point-modification} to the threefold $M$
of Theorem~\ref{thm:main} at the chosen point $x$.  If $\rho$ is its
continuous psh exhaustion, then
\[
 \rho_N=\rho\circ\mu,\qquad
 \{\rho_N\le a\}=\mu^{-1}(\{\rho\le a\})
\]
is a continuous psh exhaustion of $N$.  A holomorphic map is locally
Lipschitz in smooth coordinate charts, so local Lipschitz regularity
of the exhaustion is preserved.  A $C^2$ psh
exhaustion on $N$ would give one on $M$ by
Proposition~\ref{prop:compact-modification} with $K=\{x\}$, contrary
to Theorem~\ref{thm:main}.
\end{proof}

\begin{remark}[The sense in which the change is local]
The point $x$ can be chosen outside the zero section of $M\to S$.
The construction then preserves a neighborhood of that section and,
by properness, a neighborhood of infinity.  More precisely, for every
neighborhood $U$ of $x$, $\mu$ is an isomorphism over $M\setminus U$.
This is a modification supported over one point; it makes no claim
about a small deformation of the complex structure on a fixed smooth
manifold.
\end{remark}

\subsection{Compact products and an explicit fivefold}\label{sec:products}

\begin{proposition}[Compact product principle]\label{prop:compact-product}
Let $X$ be a nonempty complex manifold and $T$ a nonempty compact
connected complex manifold.  For each $r\in\{0,2,\infty\}$, $X$
admits a $C^r$ psh exhaustion if and only if $X\times T$ does.  If
$T$ is non-K\"ahler, then $X\times T$ is non-K\"ahler.
\end{proposition}

\begin{proof}
Pullback by the projection preserves regularity and
plurisubharmonicity, with compact sublevels
\begin{equation}\label{eq:product-sublevels}
 \{f\circ\operatorname{pr}_X\le a\}=\{f\le a\}\times T.
\end{equation}
Conversely, restriction to the closed holomorphic slice
$X\times\{t_0\}$ preserves regularity, plurisubharmonicity, and
compact sublevels.  A K\"ahler form on the product would restrict to
one on $\{x_0\}\times T$, proving the final assertion.  In fact, the
maximum principle makes every continuous psh function on the product
constant on each compact connected fiber $\{x\}\times T$.
\end{proof}

\begin{corollary}[A non-K\"ahler fivefold]\label{cor:nonkahler}
Let $M$ be the threefold of Theorem~\ref{thm:main} and let
\begin{equation}\label{eq:hopf-surface}
 H=(\C^2\setminus\{0\})/\langle z\mapsto2z\rangle.
\end{equation}
Then $M\times H$ is a smooth non-K\"ahler fivefold admitting a
continuous psh exhaustion but no $C^2$ psh exhaustion.  For every
$0<\gamma<1$ its exhaustion can be chosen locally $C^{0,\gamma}$.
\end{corollary}

\begin{proof}
The dilation action is free and properly discontinuous: if a compact
set lies in $r\le\|z\|\le R$, only finitely many integers $n$ satisfy
$r/R\le2^n\le R/r$.  Polar coordinates identify its quotient smoothly
with
\begin{equation}\label{eq:hopf-topology}
 H\simeq S^3\times\bigl(\R/(\log2)\mathbb Z\bigr)
 \simeq S^3\times S^1.
\end{equation}
Thus $H$ is a compact connected complex surface with
$H^2_{\mathrm{dR}}(H)=0$.  Any K\"ahler form $\omega$ would be exact,
and Stokes' theorem would give $\int_H\omega^2=0$, contradicting
positivity.  Apply Proposition~\ref{prop:compact-product} and
Theorem~\ref{thm:main}.  Explicitly,
\begin{equation}\label{eq:nonkahler-exhaustion}
 \rho_{M\times H}(v,[z])=\rho(v)
\end{equation}
is the required exhaustion, with the same local H\"older regularity.
\end{proof}

\section{The surface problem and boundary reductions}\label{sec:surfaces}

\subsection{Line-bundle total spaces over compact curves}
\label{sec:surface-curves}

The construction of Theorem~\ref{thm:main} has a two-dimensional compact
base and a one-dimensional fiber.  Its dimension cannot be lowered by
replacing the base by a compact Riemann surface: on a compact curve,
pseudoeffectivity of a line bundle already implies smooth semipositivity.
Here pseudoeffectivity means the existence of a closed positive
$(1,1)$-current in the first Chern class; its potentials may be singular.
Thus this assertion is stronger than the analogous assertion for a
continuous semipositive metric.

\begin{proposition}[Semipositivity on a compact curve]
\label{prop:curve-semipositivity}
Let $B$ be a compact connected Riemann surface and $A$ a holomorphic
line bundle on $B$.  The following conditions are equivalent:
\begin{enumerate}[label=\textup{(\roman*)},leftmargin=2.5em]
\item $A$ is pseudoeffective;
\item $\deg A\ge0$;
\item $A$ admits a smooth Hermitian metric with semipositive curvature.
\end{enumerate}
For $\deg A=0$ the metric can be chosen flat, and for $\deg A>0$ it can
be chosen to have strictly positive curvature.
\end{proposition}

\begin{proof}
A positive $(1,1)$-current on a curve is a positive measure.  If it
represents $c_1(A)$, its total mass is $\deg A$, proving
\textup{(i)}$\Rightarrow$\textup{(ii)}.  For the converse construction,
choose a smooth Hermitian metric $h_0$ on $A$, with curvature $\theta_0$,
and a smooth positive $(1,1)$-form $\omega$ with $\int_B\omega=1$.
Write $d=\deg A$.  Since
\[
\int_B(d\omega-\theta_0)=0,
\]
the Poisson equation, equivalently the $\partial\bar\partial$-lemma
on a compact Riemann surface, gives a smooth real function $f$ with
\[
\ddc f=d\omega-\theta_0.
\]
The metric $h=h_0e^{-f}$ has curvature
\[
\Theta_h(A)=\theta_0+\ddc f=d\omega.
\]
It is semipositive if $d\ge0$, flat if $d=0$, and strictly positive
if $d>0$.  Its curvature itself is a positive current, completing the
equivalence.
\end{proof}

\begin{remark}
This is the classical degree criterion, and does not assert
semiampleness.  A non-torsion line bundle in $\operatorname{Pic}^0(B)$
is unitary flat, but has no nonzero section of any positive tensor
power.  For a noncompact Riemann surface, every holomorphic line bundle
is holomorphically trivial: the surface is Stein, and the exponential
sequence gives this assertion from $H^1(B,\OO_B)=0$ and
$H^2(B,\mathbb Z)=0$.  The compact case is the one relevant to the
construction in this paper.
\end{remark}

\begin{corollary}[Line-bundle total spaces of dimension two]
\label{cor:curve-total-space}
For a holomorphic line bundle $A$ on a compact connected Riemann surface
$B$, the following are equivalent:
\begin{enumerate}[label=\textup{(\roman*)},leftmargin=2.5em]
\item $\Tot(A^{-1})$ admits a continuous psh exhaustion;
\item $\deg A\ge0$;
\item $\Tot(A^{-1})$ admits a smooth psh exhaustion.
\end{enumerate}
\end{corollary}

\begin{proof}
If $\deg A\ge0$, Proposition~\ref{prop:curve-semipositivity} and
Lemma~\ref{lem:norm} give the smooth squared dual norm exhaustion.
For the remaining implication, compactify the total space as
\[
\Tot(A^{-1})=
\PP_{\mathrm{lines}}(\OO_B\oplus A^{-1})\setminus C_\infty,
\qquad C_\infty=\PP_{\mathrm{lines}}(A^{-1}).
\]
Its boundary normal bundle is
\[
N_{C_\infty/\PP(\OO_B\oplus A^{-1})}
\simeq\operatorname{Hom}(A^{-1},\OO_B)\simeq A.
\]
If $\deg A<0$, this boundary is negative definite, and
Proposition~\ref{prop:surface-negative-boundary} below rules out a
continuous psh exhaustion of its complement.
\end{proof}

In particular, a continuous semipositive metric on a line bundle over
a compact curve cannot furnish the metric obstruction used in
Theorem~\ref{thm:main}: a smooth semipositive metric always exists on
that same bundle.  This rules out the present construction in complex
dimension two; it does not prove the equivalence of the two exhaustion
definitions for arbitrary complex surfaces.

\subsection{Historical surface results and their precise scope}
\label{subsec:surface-history}

The surface problem has two different aspects: regularizing an exhaustion,
and deciding when the surface is holomorphically convex. The latter
property gives a regularization mechanism through a proper holomorphic
map, but it is stronger than the existence of a smooth psh exhaustion.

Ohsawa's 1981 result \cite[Proposition~1.4]{Ohsawa1981} states that a
connected noncompact smooth weakly $1$-complete surface is
holomorphically convex if and only if it has a nonconstant holomorphic
function. His definition of weak $1$-completeness assumes a
$C^\infty$ psh exhaustion. The remark following his Theorem~1.1 says
that continuity suffices for that theorem's dichotomy concerning the
fibers of a holomorphic function; it should not be read as changing the
regularity hypothesis throughout the paper.

Tomassini and V\^aj\^aitu subsequently treated singular surfaces with a
continuous psh exhaustion \cite{TomassiniVajaitu2008}; see also the
more general holomorphic-convexity theorem
\cite[Theorem~1.1]{Vajaitu2022}. We use only its consequence for a
smooth connected surface equipped with a finite continuous psh
exhaustion. The following proposition records that consequence,
rather than a new holomorphic-convexity theorem. No assertion about
regularizing an arbitrary possibly discontinuous exhaustion is needed.

\begin{proposition}[Consequence of V\^aj\^aitu]\label{prop:surface-nonconstant-function}
Let $X$ be a connected noncompact complex surface admitting a continuous
psh exhaustion. If $\OO(X)\ne\C$, then $X$ admits a smooth psh
exhaustion.
\end{proposition}

\begin{proof}
V\^aj\^aitu's theorem makes $X$ holomorphically convex. Its Remmert
reduction is a proper holomorphic map $r:X\to Y$ to a Stein space.
Choose a smooth strongly psh exhaustion $\eta$ on $Y$, where smoothness
on a complex space means local restriction of a smooth function from a
local ambient Euclidean space. Then $\eta\circ r$ is smooth on $X$,
including along the positive-dimensional fibers, and it is psh by
holomorphic pullback. Finally,
\[
 \{\eta\circ r\le t\}=r^{-1}(\{\eta\le t\})
\]
is compact by properness. Thus $\eta\circ r$ is an exhaustion. Its
Levi form may degenerate along compact fibers, as is appropriate for
weak $1$-completeness.
\end{proof}

This implication does not dispose of surfaces for which
$\OO(X)=\C$. In particular, a holomorphically convex noncompact
connected surface always has a nonconstant holomorphic function, whereas
a smooth psh exhaustion can exist without one. Thus replacing the
holomorphic-function hypothesis in this argument requires another
regularization mechanism.


For a continuously weakly complete manifold $X$, its continuous minimal
kernel is
\[
 \Sigma_X^0=\bigcap_v\{x\in X:
       v\text{ is not strongly psh near }x\},
\]
where $v$ ranges over continuous psh exhaustions. Strong
plurisubharmonicity here means stability under sufficiently small
$C^2$ perturbations locally.

\begin{proposition}[Mongodi--Slodkowski criterion]
\label{prop:surface-minimal-kernel}
Let $X$ be a connected complex manifold with a continuous psh exhaustion
$u$. If $u(\Sigma_X^0)$ contains no terminal half-line $(a,+\infty)$,
then $X$ admits a smooth psh exhaustion.
Consequently, if a continuously weakly complete manifold has no smooth
psh exhaustion, then for every continuous psh exhaustion $u$ there is
$a_u\in\R$ such that
\[
 (a_u,+\infty)\subset u(\Sigma_X^0).
\]
\end{proposition}
\begin{proof}
By \cite[Lemma~2.1]{MongodiSlodkowski2020}, the stated image condition
makes $X$ a modification of a Stein space. Pulling back a smooth
strongly psh exhaustion by this proper holomorphic map gives the
claim. The final assertion is the contrapositive. No ambient
real-analytic exhaustion is assumed in this lemma.
\end{proof}

A separate theorem in \cite[Corollaries~4.4--4.5]{MongodiSlodkowski2020}
requires a noncompact surface $M$ already carrying a real-analytic psh
exhaustion. Every continuously weakly complete domain $\Omega\subset M$
then has a real-analytic psh exhaustion, strongly psh outside
$\Sigma_\Omega^0$, and its continuous and real-analytic minimal kernels
coincide. Their Theorem~1.1 also describes the domain's geometry through
Stein modifications, proper fibrations, or Grauert-type structure.

The ambient hypothesis needs independent verification: taking
$M=\Omega$ is circular, and a compactification
$\Omega=S\setminus|D|$ does not supply the required noncompact ambient
surface. These results do not prove the unrestricted implication
or settle every projective divisor complement. This limitation does not
assert that the unrestricted problem has any particular published status.

\subsection{Divisorial compactifications and allowed modifications}
\label{sec:surface-pair-reduction}

For the divisor-complement arguments in the rest of this section,
$S$ is a connected smooth compact complex surface, $D$ is a nonzero
effective divisor, and
\[
 X=S\setminus\supp D.
\]
The existence of this compactification is an additional hypothesis on
$X$, rather than a property of every noncompact complex surface.  Only
the support of $D$ determines the complement.  We consequently replace
$D$ by its reduced divisor whenever convenient. K\"ahlerness or
projectivity will be imposed explicitly when needed.

\begin{lemma}\label{lem:surface-boundary-divergence}
If $u$ is a continuous exhaustion of $X$, then $u(x)\to+\infty$ as
$x\to\supp D$.  More precisely, for every $b\in\R$ there is a neighborhood
$U_b$ of $\supp D$ in $S$ such that $u>b$ on $U_b\setminus\supp D$.
\end{lemma}
\begin{proof}
The sublevel set $K_b=\{u\le b\}$ is a compact subset of $X$, hence a
closed subset of $S$ disjoint from $\supp D$.  Take $U_b=S\setminus K_b$.
\end{proof}

\begin{proposition}\label{prop:surface-finite-modifications}
Let $\mu:X'\to X$ be a proper modification between connected smooth
complex surfaces with compact exceptional locus. For each
$\mathcal R\in\{C^0,C^2,C^\infty\}$, $X$ admits a psh exhaustion of
class $\mathcal R$ if and only if $X'$ does.
\end{proposition}
\begin{proof}
The image of the exceptional locus is compact, and $\mu$ is an
isomorphism outside that image. Apply
Proposition~\ref{prop:compact-modification}.
\end{proof}

\begin{remark}\label{rem:surface-boundary-transforms}
For a blowup $\mu:S'\to S$ at a boundary point, the appropriate upstairs
boundary is the full reduced inverse image of $\supp D$, including the
exceptional curve.  Its complement is biholomorphic to $X$.  In
particular, resolving the boundary to simple normal crossings does not
change the open surface.  Keeping only the strict transform of the
boundary generally changes the complement and does not give a proper
modification of the original open surface.  Proposition~
\ref{prop:surface-finite-modifications} concerns modifications with
compact exceptional locus; it does not assert descent under arbitrary
modifications with infinitely many exceptional centers escaping to
infinity.
\end{remark}

\subsection{An obstruction from negative definite boundary components}

\begin{lemma}\label{lem:surface-normal-puncture-bound}
Let $(V,p)$ be a normal complex surface germ with an isolated puncture.
A continuous plurisubharmonic function on $V\setminus\{p\}$ is bounded
above on a sufficiently small punctured neighborhood of $p$.
\end{lemma}
\begin{proof}
A local finite projection gives a proper finite holomorphic map
$\tau:V\to\Delta^2$ after shrinking representatives, with
$\tau^{-1}(0)=\{p\}$.  If $v$ is the given function, put
\[
 v_*(z)=\max\{v(y):\tau(y)=z\},\qquad z\ne0.
\]
This function is continuous, by properness and openness of the finite
map.  Away from its branch locus it is the maximum of finitely many
plurisubharmonic functions on local inverse branches.  Along the branch
locus away from $0$ it is locally bounded and continuous, so the usual
removable-singularity theorem makes it plurisubharmonic there as well.

For completeness, a Hartogs argument gives the needed bound on
$\Delta^2\setminus\{0\}$.  Choose a small bidisc and $r>0$ so that its
compact vertical circular boundary
\[
 K=\{(z_1,z_2):|z_1|\le r,\ |z_2|=r\}
\]
lies in the punctured bidisc.  Set $M=\max_K v_*$.  For $z_1\ne0$, the
maximum principle on the vertical disc gives $v_*(z_1,z_2)\le M$ for
$|z_2|<r$.  For $z_1=0$ and $z_2\ne0$, the same bound follows by
continuity.  Since $v\le v_*\circ\tau$, this bounds $v$ near $p$.
\end{proof}

\begin{proposition}\label{prop:surface-negative-boundary}
Suppose that a connected component $E$ of $\supp D$ has negative
definite intersection matrix.  Then $S\setminus\supp D$ admits no
continuous plurisubharmonic exhaustion.
\end{proposition}
\begin{proof}
Grauert's contraction theorem \cite{Grauert1962} contracts $E$, on a
neighborhood $U$ avoiding the other boundary components, to an isolated
point $p$ of a normal complex surface $V$.  The contraction induces a
biholomorphism $U\setminus E\simeq V\setminus\{p\}$.  An exhaustion
$u$ on $X$ would induce a continuous plurisubharmonic function on this
punctured neighborhood.  Lemma~\ref{lem:surface-normal-puncture-bound}
makes it bounded above near $p$, whereas Lemma~
\ref{lem:surface-boundary-divergence} requires it to tend to infinity.
\end{proof}

\subsection{Intersection matrices and smooth collars}

\begin{lemma}\label{lem:surface-collar-gluing}
Suppose that each connected component $E$ of $\supp D$ has a
neighborhood $U_E$ and a smooth plurisubharmonic function
$\Phi_E:U_E\setminus E\to\R$ tending to infinity at $E$.  Then $X$
admits a smooth plurisubharmonic exhaustion.
\end{lemma}
\begin{proof}
There are finitely many boundary components.  Choose pairwise disjoint
smaller neighborhoods $W_E$ whose closures lie in $U_E$, with compact
boundaries disjoint from $E$.  Choose $T_E>\max_{\partial W_E}\Phi_E$.
Let $\chi_E$ be smooth, convex, nondecreasing, zero on
$(-\infty,T_E]$, and unbounded at infinity.  Continuity ensures that
$\chi_E(\Phi_E)$ vanishes on a neighborhood of $\partial W_E$.
Therefore the function equal to $\chi_E(\Phi_E)$ on $W_E\setminus E$
and zero elsewhere on $X$ is smooth and plurisubharmonic.  The sum of
these functions tends to infinity at every boundary component.  Since
$S$ is compact, its sublevel sets are compact subsets of $X$.
\end{proof}

\begin{lemma}\label{lem:surface-positive-boundary-metric}
Let $E=\bigcup_{i=1}^rE_i$ be a connected reduced simple normal crossing
divisor on a smooth compact complex surface.  If a divisor
$A=\sum_i a_iE_i$, with $a_i\in\mathbb Z_{>0}$, satisfies
$A\cdot E_i>0$ for every $i$, then $\OO_S(A)$ has a smooth Hermitian
metric of strictly positive curvature on a neighborhood of $E$.
\end{lemma}
\begin{proof}
The restriction $\OO_S(A)|_E$ is ample, since its degree on every
irreducible component is positive.  Choose $m>0$ for which
$\OO_S(mA)|_E$ is very ample, and endow this restriction with its
Fubini--Study metric.  On a sufficiently small ambient coordinate
neighborhood, holomorphic sections giving the embedding of $E$ lift to
holomorphic local sections of $\OO_S(mA)$.  Their Fubini--Study weights,
divided by $m$, define local smooth metrics of $\OO_S(A)$, with the
prescribed restriction to $E$ and positive curvature in tangent
directions along its smooth part.

At a node the Zariski tangent space of $E$ equals the ambient tangent
space.  A projective embedding of $E$ is injective on that tangent
space.  Accordingly, the local Fubini--Study metric obtained by lifting
its sections has strictly positive ambient curvature at the node, and
on a smaller neighborhood of it.  Fix these local metrics on disjoint
small neighborhoods of the finitely many nodes.  Patch the remaining
local extensions of the prescribed metric using a smooth partition of
unity, relative to any global reference metric, while keeping the
chosen metrics unchanged on smaller node neighborhoods.  On the
smooth part of $E$, the tangential curvature is unchanged by this
patching: tangential second derivatives depend only on the common
restriction to the curve.  The resulting global smooth metric thus has
positive tangential curvature on $E$ and strictly positive ambient
curvature near its nodes.

Let $s_E$ be the canonical section of $\OO_S(E)$, choose any smooth
metric on this bundle, and set $q=\|s_E\|^2$.  Along the smooth part of
$E$, $\ddc q$ vanishes on the tangent line and is strictly positive in
the normal direction.  At a node $\ddc q=0$, because a reduced local
equation is the product of the two branch coordinates.  In particular,
$\ddc q$ is positive semidefinite along all of $E$.

Add $Cq$ to the local weights of the patched metric.  Away from fixed
node neighborhoods, compactness, positivity of the tangential block,
and positivity of the normal block of $\ddc q$ show that a sufficiently
large $C$ makes the full curvature matrix positive along $E$; this is
also immediate from its Schur complement.  Near the nodes, the old
curvature is already positive and $\ddc q$ is semipositive on $E$.
The new curvature is therefore positive everywhere on $E$, hence on
an ambient neighborhood by continuity.
\end{proof}

\begin{proposition}\label{prop:surface-boundary-matrix}
Assume that $S$ is K\"ahler, that $D$ is reduced with simple normal crossings and that $X$
admits a continuous plurisubharmonic exhaustion.  For each connected
component $E=\bigcup_{i=1}^rE_i$ of $D$, let
$I=(E_i\cdot E_j)$ and let $\lambda$ be its largest eigenvalue.
Then $\lambda\ge0$, and the following alternatives hold.
\begin{enumerate}[label=\textup{(\roman*)},leftmargin=2.5em]
\item If $\lambda>0$, a smooth strictly plurisubharmonic collar
exhaustion exists near $E$.
\item If $\lambda=0$, there are positive integers $a_i$ such that
$A=\sum_i a_iE_i$ is nef and
\[
 A\cdot E_i=0\quad\hbox{for every }i,\qquad A^2=0.
\]
The matrix $I$ is negative semidefinite with one-dimensional kernel.
\end{enumerate}
Thus the only connected boundary configurations not already settled by
the collar construction are supports of effective nef square-zero
divisors as in \textup{(ii)}.
\end{proposition}
\begin{proof}
The off-diagonal entries of $I$ are nonnegative, and connectedness of
$E$ makes the matrix irreducible.  Applying Perron--Frobenius to $I+c
\operatorname{Id}$ for sufficiently large $c$, its largest eigenvalue
is simple and has a strictly positive eigenvector.  If $\lambda<0$,
$I$ is negative definite, contrary to Proposition~
\ref{prop:surface-negative-boundary}.

If $\lambda>0$, approximate its positive eigenvector by a positive
rational vector $a$ sufficiently closely that $Ia$ has strictly
positive entries, and clear denominators.  Then $A\cdot E_i>0$ for
all $i$.  Lemma~\ref{lem:surface-positive-boundary-metric} gives a
positive metric $h$ on a neighborhood of $E$.  The function
\[
 \Phi=-\log\|s_A\|_h^2
\]
is smooth and strictly plurisubharmonic off $E$, since
$\ddc\Phi=\Theta_h(\OO_S(A))$ there.  It tends to infinity at every
point of $E$.

If $\lambda=0$, every eigenvalue is nonpositive, and simplicity of the
largest eigenvalue gives a one-dimensional kernel spanned by a
positive vector.  Since $I$ has integral entries, its kernel has a
rational generator; clear denominators to obtain positive integers
$a_i$.  Then $A\cdot E_i=0$ and $A^2=0$.  For any irreducible curve
$G$ outside $E$, effectiveness gives $A\cdot G\ge0$.  In the
projective case this proves nefness.  For a nonprojective compactifier,
the independent argument in Corollary~\ref{thm:surface-nonprojective-pair}
below identifies $E$ as a whole elliptic fiber support.  Its positive
kernel vector is proportional to the fiber multiplicity vector, so
some positive multiple of $A$ is a positive multiple of a fiber divisor.
Pulling back a
positive metric from the base and taking its appropriate root gives
a smooth semipositive metric on $\OO_S(A)$, proving nefness also in
the analytic sense.  Its numerical class is nonzero, because
$A\cdot H>0$ for a
K\"ahler class $H$.
\end{proof}

\subsection{Disconnected boundaries and nonprojective compactifiers}\label{sec:surface-nonprojective}

\begin{theorem}\label{thm:surface-disconnected-boundary}
Let $S$ be smooth projective and let $D$ be a nonzero effective divisor
whose support is disconnected.  If $S\setminus\supp D$ admits a
continuous plurisubharmonic exhaustion, then it admits a smooth one.
\end{theorem}
\begin{proof}
Resolve the boundary using its full reduced inverse image, which
preserves both the complement and the number of its connected boundary
components.  Write these components as $E_1,\ldots,E_r$, with $r\ge2$.
Proposition~\ref{prop:surface-boundary-matrix} supplies an effective
divisor $A_k$ supported exactly on $E_k$, with $A_k^2\ge0$; in the
positive-eigenvalue case its square is strictly positive.  Disjointness
gives $A_k\cdot A_l=0$ for $k\ne l$.

If $A_k^2>0$, the Hodge index theorem makes the intersection form
negative definite on its orthogonal complement.  The numerical class
of $A_l$ is nonzero, because its intersection with an ample divisor is
positive.  We would therefore have $A_l^2<0$, a contradiction.  Thus
every boundary component is in the zero-eigenvalue alternative:
\[
 A_k\text{ is nef},\qquad A_k^2=0,\qquad A_k\cdot A_l=0.
\]
The Hodge index theorem also implies that these nonzero isotropic nef
classes are proportional.  Their proportionality factors are positive
and rational, as can be seen by intersecting them with an integral
ample divisor.  For every $k\ge2$, choose positive integers $p_k,q_k$
such that $p_kA_k-q_kA_1$ is numerically trivial.

On a smooth projective surface, numerical triviality of a divisor
implies that its real first Chern class vanishes.  Equivalently, this
follows from the Hodge index theorem on divisor classes.  The
$\partial\bar\partial$-lemma then supplies a smooth unitary flat
Hermitian metric $h_k$ on
\[
 L_k=\OO_S(p_kA_k-q_kA_1).
\]
Indeed, the curvature of any smooth reference metric is an exact real
$(1,1)$-form, and changing its weight by a global smooth potential
makes its curvature zero.  This argument also permits torsion integral
first Chern classes.

Let $\sigma_k$ be the canonical meromorphic section of $L_k$ with
divisor $p_kA_k-q_kA_1$.  On $X=S\setminus\supp D$, the function
\[
 \phi_k=\log\|\sigma_k\|_{h_k}^2
\]
is smooth and pluriharmonic.  It tends to $-\infty$ at $E_k$ and to
$+\infty$ at $E_1$.  Hence
\[
 \rho=\sum_{k=2}^r\phi_k^2
\]
is smooth and plurisubharmonic, because
$\ddc(\phi_k^2)=\frac{i}{\pi}\partial\phi_k\wedge\barp\phi_k\ge0$.
It tends to infinity at every boundary component.  Compactness of $S$
makes it an exhaustion.  The same function transports through the
biholomorphism of complements obtained from the boundary resolution.
\end{proof}


\begin{theorem}[A numerical criterion at nonprojective infinity]
\label{thm:nonprojective}
Let $S$ be a smooth compact connected nonprojective complex surface and $D\ne0$ an effective
divisor. Put $X=S\setminus\supp D$. The following are equivalent:
\begin{enumerate}[label=\textup{(\roman*)},leftmargin=2.5em]
\item $X$ admits a continuous plurisubharmonic exhaustion;
\item $X$ admits a smooth plurisubharmonic exhaustion;
\item no connected component of $\supp D$ has negative-definite intersection matrix.
\end{enumerate}
There is no K\"ahler assumption, and the boundary need not have normal crossings.
\end{theorem}
\begin{proof}
Proposition~\ref{prop:surface-negative-boundary} proves
\textup{(i)}$\Rightarrow$\textup{(iii)}, without a K\"ahler hypothesis.
The implication \textup{(ii)}$\Rightarrow$\textup{(i)} is immediate.
We prove \textup{(iii)}$\Rightarrow$\textup{(ii)}.

The algebraic dimension of $S$ is zero or one. If $a(S)=1$, its algebraic reduction $f:S\to
B$ is a holomorphic elliptic fibration with connected fibers and every compact curve is
vertical; see \cite[Section~2]{FN2005}. For a fiber $F=\sum m_iC_i$, its intersection matrix
$I$ satisfies
\begin{equation}\label{eq:surface-fiber-intersection-identity}
 x^{\mathsf T}Ix=-\sum_{i<j}(C_i\cdot C_j)m_im_j
       \left(\frac{x_i}{m_i}-\frac{x_j}{m_j}\right)^2.
\end{equation}
Connectedness makes its kernel the line spanned by $(m_i)$. Every proper principal submatrix
is consequently negative definite. Condition \textup{(iii)} therefore forces the boundary to
contain every fiber support it meets. Hence $X=f^{-1}(B\setminus P)$ for a finite nonempty
set $P$, and the restriction of $f$ is proper. Pull back a smooth strictly subharmonic
exhaustion of $B\setminus P$.

Suppose $a(S)=0$. By \cite[Lemma 2.1(1)]{FN2005}, $c_1(L)^2\le0$ for every holomorphic line
bundle $L$ on $S$. Thus every boundary intersection matrix is negative semidefinite. On each
connected component, condition \textup{(iii)} and Perron--Frobenius give a strictly positive
kernel vector; the integral matrix permits an integral choice after rescaling. Summing the
corresponding divisors gives
\[
 A=\sum a_iD_i,\qquad a_i>0,\qquad A^2=0,
 \qquad\supp A=\supp D.
\]

Contract to a smooth minimal model $p:S\to S_0$. For any single blowdown $q:T\to T'$ with
exceptional curve $E$, put $A'=q_*A$ and write $A=q^*A'+bE$. Then
\[
 0=A^2=(A')^2-b^2.
\]
The surface $T'$ remains nonprojective, so $(A')^2\le0$ by the same lemma. Consequently
$b=0$, $(A')^2=0$, and $A'\ne0$. Iterating gives $A=p^*A_0$ for a nonzero effective
square-zero divisor $A_0$ and, crucially,
\[
 \supp A=p^{-1}(\supp A_0).
\]
Thus $X=p^{-1}(X_0)$ for $X_0=S_0\setminus\supp A_0$, and $p|_X$ is proper. A smooth psh
exhaustion of $X_0$ pulls back to one of $X$.

If $S$ is K\"ahler, its minimal model of algebraic dimension zero is a torus or a K3 surface
\cite[Section~2]{FN2005}. A torus of algebraic dimension zero has no curves: a curve on a
nonprojective torus must be a smooth elliptic curve by adjunction and the absence of
rational curves, and its elliptic quotient would give algebraic dimension at least one. On a
K3 surface, Riemann--Roch and $H^0(S_0,\OO(-A_0))=0$ yield $h^0(S_0,\OO(A_0))\ge2$,
contradicting $a(S_0)=0$. Thus this alternative cannot occur.

If $S$ is non-K\"ahler, $S_0$ is of class~VII; see \cite[Section~2]{FN2005}. For
$b_2(S_0)=n>0$, Enoki's Main Theorem \cite[p.~454]{Enoki} identifies $S_0$ with
$S_{n,\alpha,t}$ and $A_0$ with a positive multiple of its distinguished rational cycle. Its
complement is an affine-line bundle of degree $-n$ over an elliptic curve \cite[(0.4),
p.~453]{Enoki}. Completing the affine fibers gives a smooth projective $\PP^1$-bundle whose
infinity section $C_\infty$ has square $n>0$.
Lemma~\ref{lem:surface-positive-boundary-metric} gives a metric of
positive curvature near $C_\infty$. The function
$-\log\|s_{C_\infty}\|^2$ is a smooth strictly psh divergent collar,
and Lemma~\ref{lem:surface-collar-gluing} gives a smooth exhaustion of $X_0$.

For $b_2(S_0)=0$, a minimal class~VII surface with a curve is elliptic or Hopf by the
theorem of Kodaira recalled in \cite[p.~453]{Enoki}. Algebraic dimension zero excludes the
elliptic-fibration case, so $S_0$ is Hopf. Every irreducible curve $C$ on $S_0$ has
$C^2=K_{S_0}\cdot C=0$, hence arithmetic genus one. Its normalization cannot be rational:
any holomorphic map from $\PP^1$ lifts to the universal cover $\C^2\setminus\{0\}$ and is
constant. Therefore $C$ is smooth elliptic. Distinct such curves are disjoint because their
intersection numbers vanish. The boundary of $X_0$ is consequently a disjoint union of
smooth square-zero elliptic curves. Finite Ueda type gives a smooth strictly psh collar by
\cite[Theorem 2.5]{CLPT2018}; infinite type gives a unitary flat neighborhood by
\cite[Theorem 1.6]{CFSTZ2026}. The latter gives a smooth pluriharmonic logarithmic collar.
Lemma~\ref{lem:surface-collar-gluing} combines these collars into a smooth psh exhaustion of
$X_0$.
\end{proof}

\begin{remark}
The Hopf elliptic-neighborhood step uses the September 2026 preprint \cite{CFSTZ2026}.
Enoki's theorem applies to any nonzero square-zero divisor on a minimal class~VII surface
with positive second Betti number; it does not assume that the divisor is already a cycle.
Thus the proof does not use an unresolved classification of all class~VII surfaces. It also
explains why the whole pullback of the boundary, rather than its strict transform alone,
matters when passing to the minimal model.
\end{remark}


\begin{corollary}[The nonprojective K\"ahler case]
\label{thm:surface-nonprojective-pair}
Let $S$ be a connected smooth compact nonprojective K\"ahler surface
and $D\ne0$ effective. If $S\setminus\supp D$ has a continuous psh
exhaustion, then $a(S)=1$, and the algebraic reduction is a holomorphic
elliptic fibration $f:S\to B$ with connected fibers such that
\[
 \supp D=f^{-1}(P)
\]
for a finite nonempty subset $P\subset B$. A smooth exhaustion is
pulled back from $B\setminus P$.
\end{corollary}
\begin{proof}
The proof of Theorem~\ref{thm:nonprojective} excludes algebraic
dimension zero in the K\"ahler case and proves the whole-fiber
description in algebraic dimension one, including for nonminimal
surfaces and multiple fibers.
\end{proof}

\subsection{Moving boundaries and the residual projective case}

\begin{proposition}\label{prop:surface-semiample-boundary}
Let $S$ be projective and let $A$ be a nonzero effective semiample
divisor with $A^2=0$.  Then $S\setminus\supp A$ admits a smooth
plurisubharmonic exhaustion.  There are a smooth projective curve $B$,
a proper holomorphic map $f:S\to B$ with connected fibers, a positive
integer $m$, and a nonzero effective divisor $G$ on $B$ such that
\[
 mA=f^*G.
\]
\end{proposition}
\begin{proof}
A globally generated multiple of $A$ defines a projective morphism.
Its image is a curve: it is not a point because $A$ is numerically
nonzero, and it cannot have dimension two because $A^2=0$.  Its Stein
factorization gives $f:S\to B$ with $B$ a smooth projective curve and
connected fibers, and $\OO_S(mA)=f^*L$ for a line bundle $L$ on $B$
after increasing $m$ if necessary.  Since $f_*\OO_S=\OO_B$, the
canonical section with divisor $mA$ is the pullback of a section of
$L$.  Its divisor $G$ is effective and nonzero.  Consequently,
\[
 S\setminus\supp A=f^{-1}(B\setminus\supp G),
\]
and the restriction of $f$ to this complement is proper.

Every compact Riemann surface with a nonempty finite set removed has
a smooth strictly subharmonic exhaustion $\psi$.  Explicitly, fix an
area form $\omega$ of total mass one, and for each removed point $p$
take a Green function $g_p$ with $\ddc g_p=\delta_p-\omega$.  Then
$\psi=-\sum_p g_p$ is smooth with $\ddc\psi=\#\{p\}\,\omega>0$ on
the punctured curve and tends to infinity at each puncture.  The
pullback $\psi\circ f$ is the required smooth plurisubharmonic
exhaustion.
\end{proof}

\begin{proposition}\label{prop:surface-moving-boundary}
Let $S$ be smooth projective and let $A$ be a nonzero effective nef
divisor with $A^2=0$.  If either $\kappa(S,A)=1$ or $K_S\cdot A<0$,
then $S\setminus\supp A$ admits a smooth plurisubharmonic exhaustion.
\end{proposition}
\begin{proof}
A nef line bundle of Iitaka dimension one on a compact normal surface
is semiample by \cite[Lemma~7.3]{Fujino2021}.  Thus the first case
follows from Proposition~\ref{prop:surface-semiample-boundary}.

For the second case, Riemann--Roch gives
\[
 \chi(\OO_S(nA))=\chi(\OO_S)-\frac n2K_S\cdot A.
\]
For an ample divisor $H$, the intersection
$(K_S-nA)\cdot H$ is negative for large $n$, since $A\cdot H>0$.
Therefore $H^0(S,\OO_S(K_S-nA))=0$ and, by Serre duality,
$H^2(S,\OO_S(nA))=0$.  It follows that
$h^0(S,\OO_S(nA))\ge\chi(\OO_S(nA))$ is unbounded when $K_S\cdot A<0$.
Consequently $\kappa(S,A)\ge1$.  A nef divisor of square zero has
Iitaka dimension at most one: a moving linear system with
two-dimensional image would have a positive-square moving part on a
resolution, contradicting nefness and $A^2=0$.  Thus
$\kappa(S,A)=1$, and the first case applies.
\end{proof}

\begin{corollary}\label{cor:surface-connected-residual}
For divisorial compactifications by smooth projective surfaces, the
unproved continuous-to-smooth exhaustion implication reduces to the
following class of pairs:
\[
 \begin{gathered}
 D\text{ reduced with simple normal crossings and connected support},
 \\[2pt]
 A=\sum_i a_iD_i,\quad a_i\in\mathbb Z_{>0},\quad A\text{ nef},
 \\[2pt]
 A\cdot D_i=0\text{ for every }i,\qquad A^2=0,
 \\[2pt]
 \kappa(S,A)=0,\qquad K_S\cdot A\ge0.
 \end{gathered}
\]
In this reduction, the standing hypothesis that
$S\setminus\supp D$ has a continuous plurisubharmonic exhaustion is
retained.
\end{corollary}
\begin{proof}
Boundary resolution preserves the open surface.  Disconnected supports
are handled by Theorem~\ref{thm:surface-disconnected-boundary}.
Proposition~\ref{prop:surface-boundary-matrix} either gives a smooth
collar, hence a global exhaustion, or supplies $A$ with the displayed
intersection properties.  Since $A$ is nonzero and effective,
$\kappa(S,A)\ge0$; nefness and $A^2=0$ give $\kappa(S,A)\le1$.
Proposition~\ref{prop:surface-moving-boundary} settles
$\kappa(S,A)=1$ and $K_S\cdot A<0$.  The stated cases remain.
\end{proof}

\section{Birational reductions and affirmative log surface cases}
\label{sec:surface-log-cases}

The results in this section concern smooth compactifications.  Throughout,
$S$ is a smooth projective surface, $D\ne0$ is a reduced simple normal
crossing divisor, and $X=S\setminus\supp D$.  Whenever a continuous
plurisubharmonic exhaustion is assumed, it is finite-valued on $X$.
The statements do not assert that an arbitrary singular log minimal
endpoint has these smooth compactification properties.

\subsection{Birational reductions and their scope}
\label{sec:surface-mmp-scope}

The log surface MMP applies to projective pairs and, more generally,
to suitable compact K\"ahler log surface pairs; see
\cite{Fujino2012,Fujino2021}.  A nef log-adjoint endpoint is controlled
by $K_S+D$, whereas the exhaustion at infinity depends on the analytic
neighborhood of $D$.  These are distinct pieces of information.

The complement is preserved by a boundary contraction precisely when
the upstairs removed set is the full inverse image of the downstairs
removed set.  An interior contraction with compact exceptional locus
is covered by Proposition~\ref{prop:surface-finite-modifications}.
General log MMP steps need not satisfy either condition.  Moreover,
singular log endpoints may require a separate analysis of discrepancies
and boundary coefficients on a resolution.  We do not assume that
arbitrary log MMP steps preserve either exhaustion property.

For a smooth projective pair with reduced SNC boundary, write
$\overline\kappa(X)=\kappa(S,K_S+D)$ for its log Kodaira dimension.
The log MMP leads to a Mori fiber space when this dimension is
$-\infty$, and otherwise to a nef log-adjoint model.  Log abundance
then gives an associated morphism with image of dimension $0$, $1$,
or $2$.  These correspond to a numerically trivial log adjoint,
a curve Iitaka fibration, and log general type, respectively.
For a smooth curve-fibration model with $(K_S+D)\cdot F=0$, adjunction
gives $2g(F)-2+D\cdot F=0$; its general open fiber is therefore
elliptic or $\C^*$.  This organizes the cases below, but does not
assert that passage to every log endpoint preserves the open surface.
For a smooth irreducible boundary there is a useful restricted statement.
\begin{proposition}\label{prop:surface-smooth-boundary-negative-steps}
Let $S$ be smooth projective and let $C\subset S$ be a smooth connected
curve of genus at least two with $C^2=0$.  Every birational
$(K_S+C)$-negative extremal contraction from this smooth pair contracts
a $(-1)$-curve disjoint from $C$.  Thus a sequence of these birational
steps restricts to finite proper modifications of $S\setminus C$,
and both exhaustion properties are preserved.
\end{proposition}
\begin{proof}
Adjunction gives $(K_S+C)\cdot C=2g(C)-2>0$, so $C$ is not contracted.
For an irreducible contracted curve $E\ne C$, one has $E^2<0$ and
$C\cdot E\ge0$.  Again by adjunction,
\[
 (K_S+C)\cdot E=2p_a(E)-2-E^2+C\cdot E.
\]
Negativity forces $p_a(E)=0$, $E^2=-1$, and $C\cdot E=0$.
The contraction is therefore the blowdown of a smooth rational
$(-1)$-curve away from the boundary.  After each such step the surface
and the boundary remain smooth, with the same boundary genus and
self-intersection.  Proposition~\ref{prop:surface-finite-modifications}
applies at every step.
\end{proof}

\begin{remark}\label{rem:surface-log-adjoint-not-exhaustion}
Even when $K_S+D$ is semiample, its associated map need not be proper
on $S\setminus\supp D$.  Thus semiampleness of the log adjoint is not
the assertion of Proposition~\ref{prop:surface-semiample-boundary},
which requires semiampleness of a divisor supported on the boundary.
The general remaining smoothing implication is not asserted in this
article.
\end{remark}

\subsection{Ueda-theoretic sufficient conditions for a smooth boundary}

We recall the neighborhood invariant used below.  For a smooth curve
$C$ of square zero, give $N_{C/S}$ its unitary flat structure and
choose local defining functions with first-order transition constants
$t_{ij}$ of modulus one.  A system of type $n$ satisfies
\[
 w_i=t_{ij}w_j+O(w_j^{n+1}).
\]
The Ueda type is the largest attainable $n$, or infinity if such systems
exist at every order.  The successive obstructions lie in
$H^1(C,N_{C/S}^{-n})$.  Infinite type is a formal condition and does
not by itself supply convergent defining functions.  Exact equations
$w_i=t_{ij}w_j$ on a neighborhood give type $(\beta)$; infinite type
without such equations is type $(\gamma)$
\cite[\S\S2,5]{Ueda1983}.
The Diophantine condition on $N$ means that, for some $c,\tau>0$ and
an invariant distance on $\operatorname{Pic}^0(C)$,
\[
 \operatorname{dist}(N^{\otimes m},\OO_C)\ge c m^{-\tau}
 \qquad(m\ge1).
\]

\begin{proposition}[A smooth square-zero boundary]
\label{prop:surface-ueda-sufficient}
Let $C$ be a smooth connected curve on a smooth compact K\"ahler
surface $S$, with $C^2=0$, and set $N=\OO_C(C)$.  The complement
$S\setminus C$ admits a smooth psh exhaustion in each of the
following cases:
\begin{enumerate}[label=\textup{(\roman*)},leftmargin=2.5em]
 \item the Ueda type of $C\subset S$ is finite;
 \item $\OO_S(C)$ is unitary flat on a neighborhood of $C$;
 \item $N$ is torsion;
 \item for some $k>0$, the unitary representation associated with
 $N^{\otimes k}$ extends from $\pi_1(C)$ to $\pi_1(S)$;
 \item $N$ satisfies the Diophantine condition;
 \item $C$ is elliptic.
\end{enumerate}
\end{proposition}

\begin{proof}
Case \textup{(i)} follows from Ueda's strictly psh collar theorem
\cite[Theorem~2.5]{CLPT2018} and
Lemma~\ref{lem:surface-collar-gluing}.  In case \textup{(ii)},
unitary defining functions give the pluriharmonic collar function
$-\log|w_i|^2$, to which the same lemma applies.

For \textup{(iii)}, finite Ueda type has already been treated.  In
infinite type, the fibration theorem for torsion normal bundle
\cite[Theorem~1.3]{CLPT2018} makes a positive multiple of $C$ a
fiber of a holomorphic fibration; a local base coordinate gives the
required collar function.

For \textup{(iv)}, torsion normal bundles have already been treated
in \textup{(iii)}, so only infinite Ueda type with non-torsion normal
bundle remains.  Neeman's
extension theorem \cite[Theorem~4.4]{CLPT2018} gives an effective
divisor $E$ disjoint from $C$ with $[E]=m[C]$ in real cohomology for
some $m>0$. Hence $\OO_S(mC-E)$ has zero real first Chern class and
admits a unitary flat metric $h$ by the $\partial\bar\partial$-lemma.
If $\sigma$ is its canonical meromorphic section, then
$-m^{-1}\log\|\sigma\|_h^2$ is a smooth pluriharmonic collar diverging
at $C$. This also gives the unitary flat neighborhood consequence.  In \textup{(v)}, finite
type is already covered, and Ueda's convergence theorem
\cite[Theorem~3]{Ueda1983} gives the flat neighborhood in infinite type.
Finally, finite type again handles one alternative in \textup{(vi)};
for infinite type,
\cite[Theorem~1.6]{CFSTZ2026} excludes type $(\gamma)$ and hence
gives case \textup{(ii)}.
\end{proof}


\subsection{Homology and irregularity of a square-zero boundary}

\begin{lemma}[Exact character-extension criterion]\label{lem:characters}
Let $i:C\hookrightarrow S$ be the inclusion of a smooth connected compact complex curve in a
compact complex surface, and let
\[
 j=i_*:H_1(C,\mathbb Z)\longrightarrow H_1(S,\mathbb Z).
\]
A character $\rho:H_1(C,\mathbb Z)\to\Uone$ extends to $H_1(S,\mathbb Z)$ if and only if it
is trivial on $\ker j$. A positive power of $\rho$ extends if and only if $\rho(\ker j)$ is
finite. In particular, injectivity of $j$ implies extension of every unitary character,
without taking a power.
\end{lemma}
\begin{proof}
Necessity is immediate. If $\rho$ is trivial on $\ker j$, define a character on the subgroup
$j(H_1(C,\mathbb Z))$ by $j(a)\mapsto\rho(a)$. Every character from a subgroup of a finitely
generated abelian group to $\Uone$ extends to the whole group. Indeed, $\Uone$ is divisible,
and hence an injective abelian group. Alternatively one can use the structure theorem for
finitely generated abelian groups and successively take roots in $\Uone$. This argument
allows torsion in $H_1(S,\mathbb Z)$ and a nonprimitive image of $j$.

Apply the first assertion to $\rho^k$. Such a power extends exactly when $\rho(\ker j)$ is
killed by a positive integer $k$, equivalently is contained in the finite group of $k$th
roots of unity. Conversely every finite subgroup of $\Uone$ is killed by its order. Finally,
representations into $\Uone$ factor through the abelianization, so the same statements hold
for extension from $\pi_1(C)$ to $\pi_1(S)$.
\end{proof}

\begin{theorem}[A homological sufficient condition]\label{thm:homology}
Let $S$ be a smooth compact connected K\"ahler surface and $C\subset S$ a smooth connected
curve with $C^2=0$. Let $\rho$ be the unitary character associated with $N_{C/S}$, viewed as
a character of $H_1(C,\mathbb Z)$. If
\[
 \rho\bigl(\ker(H_1(C,\mathbb Z)\to H_1(S,\mathbb Z))\bigr)
 \quad\text{is finite},
\]
then $S\setminus C$ admits a smooth plurisubharmonic exhaustion. In particular this holds
when $H_1(C,\mathbb Z)\to H_1(S,\mathbb Z)$ is injective. No continuous exhaustion is
assumed.
\end{theorem}
\begin{proof}
The degree-zero normal bundle is unitary flat. By
Lemma~\ref{lem:characters}, a positive power of its character extends
to $H_1(S,\mathbb Z)$ and hence to $\pi_1(S)$. Apply
Proposition~\ref{prop:surface-ueda-sufficient}\textup{(iv)}.
\end{proof}

\begin{corollary}[Irregularity at least the boundary genus]\label{cor:irregularity}
Let $S$ and $C$ be as in Theorem~\ref{thm:homology}. If
\[
 q(S):=h^1(S,\OO_S)\ \ge\ g(C),
\]
then $S\setminus C$ admits a smooth plurisubharmonic exhaustion. No continuous exhaustion is assumed.
\end{corollary}
\begin{proof}
Consider the restriction map
\[
 r:H^1(S,\OO_S)\longrightarrow H^1(C,\OO_C).
\]
If $r$ is not injective, \cite[Proposition 2.7]{CLPT2018} implies that $C$ is a multiple
fiber of a holomorphic fibration $f:S\to B$. Taking connected fibers, $f^{-1}(b)$ has
support exactly $C$ for one point $b\in B$. Thus the restriction $S\setminus C\to
B\setminus\{b\}$ is proper; pull back a smooth strictly subharmonic exhaustion of the
punctured compact curve.

If $r$ is injective, the inequality $q(S)\ge g(C)$ forces equality, and $r$ is an
isomorphism. The functorial Hodge decomposition on compact K\"ahler manifolds identifies $r$
with the $(0,1)$ part of $i^*:H^1(S,\C)\to H^1(C,\C)$; complex conjugation identifies its
$(1,0)$ part. Therefore $i^*$ is an isomorphism, and so is $i_*:H_1(C,\mathbb Q)\to
H_1(S,\mathbb Q)$.

The kernel of the integral map $j:H_1(C,\mathbb Z)\to H_1(S,\mathbb Z)$ is consequently
torsion. Since $H_1(C,\mathbb Z)\simeq\mathbb Z^{2g(C)}$ is torsion-free, that kernel is
zero. Torsion in the target causes no difficulty: an element mapping to torsion maps to zero
rationally, and hence must already be zero in the source. Apply Theorem~\ref{thm:homology};
Lemma~\ref{lem:characters} shows that here the character itself extends, not merely a
positive power.
\end{proof}

\begin{remark}[The genus-zero endpoint]
If $C\simeq\PP^1$ and $C^2=0$, then $N_{C/S}\simeq\OO_{\PP^1}$. All Ueda obstruction groups
$H^1(C,N_{C/S}^{-n})$ vanish, so the type is infinite. The torsion-normal-bundle fibration
theorem gives a fibration with fiber $C$, and hence a smooth exhaustion of the complement.
This is also covered by Corollary~\ref{cor:irregularity}, since $q(S)\ge0=g(C)$.
\end{remark}

\subsection{Smooth rulings in every boundary degree}

In a smooth ruling the square-zero alternative has a particularly
rigid form. The following proposition includes affine-line complements
and bisections, and removes any degree restriction.

\begin{proposition}[Connected square-zero boundaries in a smooth ruling]
\label{prop:ruled-all-degrees-zero}
Let $f:S\to B$ be a smooth $\PP^1$-bundle over a connected smooth
projective complex curve. Let $E=\bigcup_{i=1}^r E_i$ be a nonempty connected
reduced simple normal crossing divisor. Suppose that there are integers
$a_i>0$ such that
\[
 A=\sum_{i=1}^r a_iE_i,\qquad A\cdot E_i=0\quad(1\le i\le r).
\]
Then $S\setminus E$ admits a smooth plurisubharmonic exhaustion, without
assuming a continuous exhaustion. More precisely, either $E$ is one fiber,
or $E=C$ is a single smooth curve with $C^2=0$ and $f|_C:C\to B$ finite
\'etale. If its degree $d$ is at least three, a finite \'etale base change
identifies the pair with
\[
 (S',E')\simeq
 \bigl(B'\times\PP^1,\ B'\times Q\bigr),
 \qquad |Q|=d.
\]
\end{proposition}

\begin{proof}
Every vertical irreducible curve in $S$ is a whole fiber, since the fibers
are irreducible. If all components of $E$ are vertical, their disjointness
and connectedness of $E$ make $E$ one fiber. The map
$S\setminus E\to B\setminus\{b\}$ is proper, so the pullback of a smooth
subharmonic exhaustion of the punctured compact curve is an exhaustion.
Such an exhaustion can, for example, be obtained from a metric of positive
curvature on $\mathcal O_B(b)$ and the negative logarithm of the squared
norm of its canonical section, followed by adding a constant.

Now suppose that $E$ has a horizontal component. For a fiber class $F$,
we have $A\cdot F>0$. Thus no fiber can be a component of $E$, since
$A\cdot E_i=0$. All components are therefore horizontal.
A $\PP^1$-bundle over a smooth projective complex curve is the
projectivization of a rank-two vector bundle; in particular $S$ is
projective and has a section $C_0$. Write
\[
 C_0^2=-e,\qquad C_0\cdot F=1,\qquad F^2=0.
\]
These classes form a basis of the real numerical divisor space.
Put $a=A\cdot F>0$ and write $A\equiv aC_0+bF$.
The orthogonality assumptions give $A^2=\sum_i a_i(A\cdot E_i)=0$,
so $b=ae/2$. Write $E_i\equiv d_iC_0+b_iF$, where
$d_i=E_i\cdot F>0$. Then
\[
 0=A\cdot E_i=a\bigl(b_i-d_ie/2\bigr),
 \qquad E_i\equiv d_i\bigl(C_0+(e/2)F\bigr).
\]
It follows that
\[
 E_i^2=0,\qquad E_i\cdot E_j=0\quad(i\ne j).
\]
Distinct irreducible curves on a smooth surface have nonnegative local
intersection multiplicities. Hence the $E_i$ are pairwise disjoint;
connectedness implies that $E=C$ is one component. The SNC hypothesis
makes $C$ smooth, and the preceding calculation gives $C^2=0$.

Set $d=C\cdot F$. The canonical numerical class is
\[
 K_S\equiv-2C_0+(2g(B)-2-e)F,
 \qquad C\equiv d\bigl(C_0+(e/2)F\bigr).
\]
Consequently
\[
 K_S\cdot C=d(2g(B)-2),\qquad
 2g(C)-2=d(2g(B)-2).
\]
The second equality is adjunction. Since the nonconstant map
$f|_C:C\to B$ is finite of degree $d$, Riemann--Hurwitz makes its
ramification divisor zero. It is therefore \'etale for every $d$.

If $d=1$, the complement is a holomorphic affine-line bundle over $B$.
Its associated linear bundle has degree zero, because its dual is the
normal bundle of the infinity section $C$, whose degree is $C^2=0$.
We recall the affine-bundle argument, including the normal form needed
for smoothness. Give the linear bundle a unitary flat metric and choose
local affine coordinates with transitions
$w_i=a_{ij}w_j+b_{ij}(z)$, where the $a_{ij}$ are unitary constants.
Choose a smooth section with local coordinates $s_i$. The forms
$\bar\partial s_i$ give a global linear-bundle-valued $(0,1)$-form.
Subtracting a global smooth section of the linear bundle makes this form
harmonic. In a local flat frame over a curve, its coefficient is
antiholomorphic. On sufficiently small coordinate disks choose
antiholomorphic primitives $g_i$ of this harmonic form. Then $s_i-g_i$
is holomorphic, and replacing $w_i$ by $w_i-(s_i-g_i)$ gives transitions
\[
 w_i=a_{ij}w_j+c_{ij},\qquad
 g_i=a_{ij}g_j+c_{ij},\qquad |a_{ij}|=1,
\]
with constant $c_{ij}$: the differences are both holomorphic and
antiholomorphic. This is the curve case of the normal form in
\cite[Proposition~2.1]{KoikeUeda2025}.
The functions $\rho_i=|w_i-g_i|^2$ agree globally. In coordinates $(z,w_i)$,
\[
 (\rho_i)_{z\bar z}=|(g_i)_{\bar z}|^2,\qquad
 (\rho_i)_{w_i\bar w_i}=1,\qquad
 (\rho_i)_{z\bar w_i}=(\rho_i)_{w_i\bar z}=0.
\]
Thus $\rho$ is smooth and plurisubharmonic. Its sublevel sets are bundles
of closed disks of bounded radius around a smooth section over compact
$B$, and hence are compact.

For $d\ge2$, choose a connected finite \'etale Galois cover
$p:B'\to B$ killing the finite monodromy of $C\to B$.
The pullback $S'=S\times_B B'$ has boundary $E'$ consisting of $d$
pairwise disjoint sections. If $q:S'\to S$ is the induced map and $T$
is any of these sections, the projection formula gives
\[
 T^2=(q^*C)\cdot T=C\cdot q_*T=0,
\]
since $q^*C$ is the reduced union of these disjoint sections and $q_*T$
is a positive integral multiple of $C$.

If $d=2$, putting the sections at zero and infinity gives transition
functions $w_i=\lambda_{ij}(z)w_j$. Thus
$S'\setminus E'\simeq L^\times$ for a line bundle $L$ of degree zero;
the degree is, up to sign, the self-intersection of either section.
For a unitary flat metric on $L$, the function
$t=\log\|v\|^2$ is pluriharmonic on $L^\times$, and $t^2$ is a smooth
nonnegative plurisubharmonic exhaustion.

If $d\ge3$, use the unique fiberwise projective transformation putting
three labeled sections at $0,1,\infty$. It varies holomorphically on
local trivializations; uniqueness makes the resulting trivializations
agree. Hence $S'\simeq B'\times\PP^1$. Every remaining section
is a holomorphic map $B'\to\PP^1$ avoiding $0,1,\infty$. Such a
map is constant, since a nonconstant map between compact connected
curves is surjective. Thus $E'=B'\times Q$ for a fixed set of $d$ points.
Writing $Q=\{\infty,q_1,\ldots,q_{d-1}\}$, an explicit smooth
nonnegative subharmonic exhaustion of $\PP^1\setminus Q$ is
\[
 r(z)=|z|^2+\sum_{j=1}^{d-1}|z-q_j|^{-2}.
\]
Its pullback to $B'\times(\PP^1\setminus Q)$ is proper because
$B'$ is compact.

In either case $d\ge2$, let $\rho'$ be the nonnegative exhaustion just
constructed and let $G$ be the deck group of $S'\setminus E'\to S\setminus C$.
The average
\[
 \rho'_G=\frac1{|G|}\sum_{g\in G}\rho'\circ g
\]
is smooth, plurisubharmonic, and invariant. It remains proper, because
$\rho'_G\ge\rho'/|G|$. It descends smoothly through the finite \'etale
cover to the required exhaustion of $S\setminus C$.
\end{proof}

\begin{corollary}[The boundary-matrix criterion for smooth rulings]
\label{cor:ruled-all-degrees}
Let $f:S\to B$ be a smooth $\PP^1$-bundle over a connected smooth
projective complex curve, and let $D\ne0$ be a reduced SNC divisor.
The following conditions are equivalent:
\begin{enumerate}[label=\textup{(\roman*)},leftmargin=2.5em]
\item $S\setminus D$ admits a continuous plurisubharmonic exhaustion;
\item $S\setminus D$ admits a smooth plurisubharmonic exhaustion;
\item The intersection matrix of every connected component of $D$ is
not negative definite.
\end{enumerate}
In particular, there is no restriction on $D\cdot F$ or on the number
of connected boundary components.
\end{corollary}

\begin{proof}
Smooth implies continuous, and
Proposition~\ref{prop:surface-negative-boundary} gives the necessary
matrix condition. Conversely, consider one connected component
$E=\bigcup_i E_i$ of $D$ and its intersection matrix $I$. Its largest
eigenvalue $\lambda$ is nonnegative by hypothesis. The
Perron--Frobenius argument in the proof of
Proposition~\ref{prop:surface-boundary-matrix} uses only this matrix
condition for the remaining conclusions. If $\lambda>0$, it gives a
positive integral divisor $H$ supported on $E$ with $H\cdot E_i>0$,
and hence a smooth strictly psh divergent collar by
Lemma~\ref{lem:surface-positive-boundary-metric}. If $\lambda=0$,
the one-dimensional rational kernel has a strictly positive integral
generator. Proposition~\ref{prop:ruled-all-degrees-zero} then supplies
a smooth exhaustion of $S\setminus E$, whose restriction gives a
collar at $E$. Lemma~\ref{lem:surface-collar-gluing} combines the
collars at all connected boundary components.
\end{proof}

\begin{remark}[Scope]
This criterion uses the rank-two numerical divisor space and the absence
of reducible fibers of a smooth ruling. It does not settle a ruled surface
after arbitrary blowups, a singular rational fibration, or the general
projective square-zero boundary problem.
\end{remark}

\subsection{Log Fano and log Calabi--Yau pairs}

\subsubsection{The log Fano case}

\begin{proposition}[The smooth log Fano case]
\label{prop:surface-log-fano}
Suppose that $-(K_S+D)$ is ample.  If $X$ admits a continuous psh
exhaustion, then it admits a smooth psh exhaustion.
\end{proposition}

\begin{proof}
Apply Proposition~\ref{prop:surface-boundary-matrix} to each connected
component $Y$ of $\supp D$.  The negative definite case is excluded
by the continuous exhaustion hypothesis.  A positive largest eigenvalue
provides a smooth strictly psh collar function, which is sufficient by
Lemma~\ref{lem:surface-collar-gluing}.  In the remaining case there
is an effective nef integral divisor $A_Y$, supported on $Y$, with
$A_Y^2=0$ and $A_Y\cdot D_i=0$ for every component of $Y$.
The other connected boundary components are disjoint from $Y$, so
$D\cdot A_Y=0$.  Therefore
\[
 K_S\cdot A_Y=(K_S+D)\cdot A_Y<0.
\]
Proposition~\ref{prop:surface-moving-boundary} gives a smooth psh
exhaustion of $S\setminus Y$, whose restriction is a collar function
at $Y$.  Choose disjoint collars for the finitely many connected
boundary components and apply the construction of
Lemma~\ref{lem:surface-collar-gluing} separately on each collar.
The sum of these functions is a smooth psh exhaustion of $X$.
\end{proof}

\subsubsection{The log Calabi--Yau case}

\begin{proposition}[Connected reduced boundary in the log Calabi--Yau case]
\label{prop:surface-log-cy}
Assume that $D$ is connected and that $K_S+D\equiv0$ numerically.
If $X$ admits a continuous psh exhaustion, then it admits a smooth
psh exhaustion.
\end{proposition}

\begin{proof}
For every irreducible component $D_i$, adjunction gives
\begin{equation}\label{eq:surface-cy-adjunction}
 0=(K_S+D)\cdot D_i
   =2g(D_i)-2+(D-D_i)\cdot D_i.
\end{equation}
Since $D$ is reduced SNC and connected, this implies that $D$ is either
a single smooth elliptic curve or a cycle of smooth rational curves.
Here a cycle with two components has two transverse intersection points.

The negative definite and positive largest-eigenvalue cases of the
boundary intersection matrix are treated by
Proposition~\ref{prop:surface-boundary-matrix}.  It remains to consider
the negative semidefinite case with a strictly positive kernel vector.
For a single elliptic component this means $D^2=0$.
If its Ueda type is finite, Ueda's theorem provides a smooth strictly
psh collar function tending to infinity at $D$
\cite[Theorem~2.5]{CLPT2018}.  If its Ueda type is infinite, the
exclusion of type $(\gamma)$ for elliptic curves in compact complex
surfaces \cite[Theorem~1.6]{CFSTZ2026} gives a unitary flat
neighborhood.  In the latter neighborhood, unitary defining functions
$w_i=t_{ij}w_j$, $|t_{ij}|=1$, give the smooth pluriharmonic function
$-\log|w_i|^2$.  Lemma~\ref{lem:surface-collar-gluing} concludes
the elliptic case.

Now suppose that $D$ is a rational cycle.  Contract any component
which is a $(-1)$-curve.  This is the inverse of blowing up a boundary
node and taking the full reduced total transform.  If $q:S\to S'$
is this contraction and $D'$ is the image cycle, then
\[
 q^*D'=D+E,\qquad K_S=q^*K_{S'}+E,
 \qquad K_S+D=q^*(K_{S'}+D').
\]
Thus $S\setminus D\simeq S'\setminus D'$ and
$K_{S'}+D'\equiv0$.  Contracting a component of a three-component
cycle yields a two-component cycle; contracting a component of a
two-component cycle yields an irreducible nodal rational curve.
Every intermediate support is a reduced nodal cycle.

The existence of a positive square-zero weighted divisor is also
preserved.  Indeed, a nef kernel divisor has intersection zero with
the contracted curve and hence is the numerical pullback of its
pushforward.  Stop when the boundary has no smooth $(-1)$-component,
or when it has become irreducible nodal.  In the irreducible case,
the existence of a nonzero square-zero weighted divisor implies
$D^2=0$.

In the remaining reducible case write $b_i=-D_i^2$ and let
$A=\sum a_iD_i$, $a_i>0$, be a kernel divisor.  Negative
semidefiniteness, positive intersections with the adjacent components,
and the absence of $(-1)$-components imply $b_i\ge2$.
Adding the equations $A\cdot D_i=0$ gives
\[
 \sum_i (b_i-2)a_i=0,
\]
where the two edges incident to a component are counted separately
also for a two-component cycle.  Hence $b_i=2$ for all $i$.
The kernel equations on a connected cycle then force all $a_i$ to
be equal.  Consequently the reduced divisor $D$ itself is nef and
$D^2=0$.

The resulting surface is rational.  Indeed, $K_S\equiv-D$ implies
$K_S\cdot H<0$ for an ample divisor $H$, so $\kappa(S)=-\infty$.
If $S$ were nonrational ruled, its rational curves would be vertical
for the ruling over a curve of positive genus.  The supports of fibers
of such a ruling, including after point blowups, have tree dual graphs;
they contain neither a rational cycle nor an irreducible nodal rational
curve.  The Enriques--Kodaira classification therefore makes $S$
rational.  On a rational surface numerical triviality of a line bundle
implies its triviality, so
\[
 \OO_S(D)\simeq K_S^{-1}.
\]

The normal line bundle $N_D=\OO_D(D)$ has multidegree zero.  Since
$D$ is a rational nodal cycle, $\operatorname{Pic}^0(D)\simeq\C^*$.
Let $q_D\in\C^*$ be its period.  Changing the orientation of the
cycle replaces $q_D$ by $q_D^{-1}$, so the following alternatives do
not depend on that choice.

If $|q_D|=1$, the semipositivity theorem for generalized Halphen
surfaces gives a smooth semipositive Hermitian metric $h$ on
$\OO_S(D)$ \cite[Proposition~6.2\textup{(B)}]{CFSTZ2026}.
The function $-\log\|s_D\|_h^2$ on $S\setminus D$ is therefore
a smooth psh exhaustion.

If $|q_D|\ne1$, then every $N_D^{-m}$ is a nontrivial multidegree-zero
line bundle on the rational cycle.  The normalization exact sequence
gives
\[
 H^0(D,N_D^{-m})=H^1(D,N_D^{-m})=0\qquad(m\ge1).
\]
In particular, all the nodal Ueda obstruction classes vanish, so the
pair has infinite type.  Koike's theorem for nonunitary normal periods
\cite[Theorem~1.6]{KoikeNodes2017} then supplies a smooth strictly
psh function on a punctured neighborhood, tending to infinity at $D$.
Lemma~\ref{lem:surface-collar-gluing} completes this last case.
\end{proof}

\begin{remark}\label{rem:surface-cy-recent}
The elliptic infinite-type step and the unitary-period rational-cycle
step above use the preprint \cite{CFSTZ2026}, submitted in September
2026.  They are stated here with that status.  The proof does not
identify finite Ueda type with smooth semipositivity of $\OO_S(D)$:
in that case it uses the collar function supplied by Ueda's theorem.
\end{remark}

\subsection{Elliptic fibrations and vertical boundaries}

\begin{proposition}[The vertical-boundary case]
\label{prop:surface-elliptic-iitaka}
Let $f:S\to B$ be a morphism with connected fibers onto a smooth
projective curve.  Suppose that its general fiber $F$ is a smooth
elliptic curve and $(K_S+D)\cdot F=0$.  If $X$ admits a continuous
psh exhaustion, then there is a finite nonempty set $P\subset B$ such
that
\[
 \supp D=f^{-1}(P),\qquad X=f^{-1}(B\setminus P).
\]
In particular, $X$ admits a smooth psh exhaustion.  This applies when
$f$ is the log Iitaka fibration with elliptic general fiber.
\end{proposition}

\begin{proof}
Since $F^2=0$ and $g(F)=1$, adjunction gives $K_S\cdot F=0$.
Thus $D\cdot F=0$, and every irreducible component of $D$ is
vertical.  Only finitely many fibers meet the boundary.

For a fiber $F_b=\sum_i m_iC_i$, the intersection identity
\eqref{eq:surface-fiber-intersection-identity} shows that every proper
principal submatrix is negative definite: its vector extends by zero
on an omitted component and therefore cannot lie in the kernel spanned
by the positive multiplicity vector.

If the boundary met $F_b$ without containing its whole support, each
of its connected boundary subconfigurations in that fiber would have
negative definite intersection matrix.  Such a connected component
is prohibited by Proposition~\ref{prop:surface-boundary-matrix}.
Consequently the boundary contains whole fiber supports.  Multiple
fibers cause no change, since the conclusion concerns supports.
The restricted map $f|_X:X\to B\setminus P$ is proper.  Pull back a
smooth strictly subharmonic exhaustion of the punctured base.
\end{proof}

\section{Examples and remaining boundary geometry}\label{sec:surface-examples}

\subsection{Smooth exhaustions without holomorphic convexity}

The log general type condition does not force holomorphic convexity of
the complement.  The following affirmative examples also show that
ramification of the boundary over an Iitaka base is compatible with a
smooth exhaustion.

\begin{proposition}\label{prop:surface-branched-flat-examples}
Let $B$ be a smooth projective curve of genus at least two, and let
$L\in\operatorname{Pic}^0(B)$ be non-torsion.  Put
\[
 S_0=\PP_{\mathrm{lines}}(\OO_B\oplus L),\qquad
 \pi:S_0\to B,\qquad C_\infty=\PP_{\mathrm{lines}}(L).
\]
For an integer $n\ge1$, choose a line bundle $A$ of sufficiently large
degree that the divisor class
$H_n=nC_\infty+\pi^*A$ is very ample.  Choose a smooth branch divisor
$R\in|2H_n|$ transverse to $C_\infty$, let $q:T\to S_0$ be its double
cover, and put $D=q^{-1}(C_\infty)$.  Then $T$ is smooth projective,
$D$ is smooth connected with $D^2=0$, and $X=T\setminus D$ admits a
smooth psh exhaustion but satisfies $\OO(X)=\C$.  Moreover,
\[
 \kappa(T,D)=0,\qquad
 \kappa(T,K_T+D)=
 \begin{cases}1,&n=1,\\2,&n\ge2.\end{cases}
\]
For $n\ge3$, $K_T$ itself is ample.
\end{proposition}

\begin{proof}
The curve $C_\infty$ has normal bundle $L^{-1}$ and square zero; it is
nef.  The ruled-surface canonical bundle formula, with our lines
convention, is
\[
 K_{S_0}\sim-2C_\infty+\pi^*(K_B\otimes L^{-1}).
\]
The double cover is defined by the algebra
$\OO_{S_0}\oplus\OO_{S_0}(-H_n)$ and the chosen section cutting out
$R$.  Smoothness of $R$ makes $T$ smooth, and transversality makes
$D$ smooth.  The restriction $D\to C_\infty$ is a double cover
branched at $R\cap C_\infty$, which has
$2\deg A>0$ points.  This implies connectedness of $D$.
The projection formula gives $D^2=2C_\infty^2=0$.

The double-cover canonical bundle formula gives
\begin{align*}
 K_T+D&\sim q^*\bigl((n-1)C_\infty+
       \pi^*(K_B\otimes L^{-1}\otimes A)\bigr),\\
 K_T&\sim q^*\bigl((n-2)C_\infty+
       \pi^*(K_B\otimes L^{-1}\otimes A)\bigr).
\end{align*}
The base line bundle in these formulas has positive degree.  For
$n=1$, the log adjoint is pulled back from this ample base bundle and
has Iitaka dimension one.  For $n\ge2$, the first divisor is ample:
it has positive intersection with every curve and positive square,
by nefness of $C_\infty$ and the ruled-surface intersection formulas.
Nakai--Moishezon and finite pullback give ampleness of $K_T+D$.
The same argument applies to $K_T$ for $n\ge3$.

There is an identification $S_0\setminus C_\infty=\Tot(L)$.
Equip $L$ with its unitary flat metric.  The squared norm
$\rho_0(v)=\|v\|^2$ is a smooth psh exhaustion of this total space.
The finite map $q|_X$ is proper, so
\[
 \rho=\rho_0\circ q|_X
\]
is a smooth psh exhaustion of $X$.

To check the absence of holomorphic functions, expand a holomorphic
function on $\Tot(L)$ in fiber Taylor series.  The coefficient of
degree $j$ is a global section of $L^{-j}$.  For $j>0$ this bundle is
nontrivial of degree zero and has no section.  The degree-zero
coefficient is constant on $B$.  Therefore
$\OO(\Tot(L))=\C$.  If $f\in\OO(X)$, its trace and norm under the
finite double cover are holomorphic functions on $\Tot(L)$, hence
constants $a$ and $b$.  Thus $f^2-af+b=0$.  Connectedness of $X$
forces $f$ to be constant.

Finally, the ratio of any section of $\OO_T(mD)$ to the canonical
section $s_D^m$ is holomorphic on $X$.  The preceding argument implies
that it is constant.  Hence $h^0(T,\OO_T(mD))=1$ for all $m\ge1$,
and $\kappa(T,D)=0$.
\end{proof}

\begin{remark}
For $n=1$, a general fiber of $T\to B$ is a double cover of
$\PP^1$ branched at two points, hence is rational; removing $D$
gives $\C^*$.  Nevertheless, $D\to B$ is ramified at
$R\cap C_\infty$.  Thus a proof of the general rational-fiber
Iitaka case cannot assume that its boundary multisection is
unramified merely because a smooth psh exhaustion exists.
For $n\ge3$ these examples have an ample ordinary canonical bundle
as well as an ample log adjoint, while their complements are not
holomorphically convex.  They are examples with smooth exhaustions,
not counterexamples to any continuous-to-smooth implication.
\end{remark}

\subsection{Square-zero boundaries with prescribed normal bundle}

The preceding propositions do not prove the smoothing implication
for every log Iitaka pair with rational general fiber, nor for every
log general type pair.  The following constructions explain why log
Kodaira dimension alone does not make the residual boundary geometry
finite or automatic.  They are numerical test families, not
counterexamples to the smoothing implication; no continuous exhaustion
is asserted for arbitrary choices of their parameters.

\begin{proposition}[Square-zero boundaries with arbitrary normal bundle]
\label{prop:surface-arbitrary-normal}
Let $S_0=\PP^1\times\PP^1$, and let $C\subset S_0$ be a smooth
curve of bidegree $(a,b)$, where $a,b\ge2$.  For every
$N\in\operatorname{Pic}^0(C)$ there are $2ab$ distinct points on
$C$ such that, after blowing them up and taking the strict transform
$D$ of $C$, one has
\[
 \begin{gathered}
 D^2=0,\qquad N_{D/S}\simeq N,\qquad D\text{ is nef},\\
 \OO_S(K_S+D)\simeq\pi^*\OO_{S_0}(a-2,b-2).
 \end{gathered}
\]
If $N$ is non-torsion, then $\kappa(S,D)=0$.
If $a=2$, $b\ge3$, the pair is log minimal of log Kodaira dimension
one.  If $a,b\ge3$, it is log minimal of log Kodaira dimension two.
\end{proposition}

\begin{proof}
The genus and self-intersection are
\[
 g(C)=(a-1)(b-1),\qquad C^2=2ab.
\]
The line bundle $\OO_C(C)\otimes N^{-1}$ has degree $2ab$, which
exceeds $2g(C)$.  It is therefore basepoint-free and a general
section has $2ab$ distinct zeros.  Let $Z$ be this reduced divisor
of zeros and let $\pi:S\to S_0$ be the blowup of its points.
The normal-bundle formula for a strict transform gives
\[
 N_{D/S}\simeq\OO_C(C-Z)\simeq N,
\]
and $D^2=C^2-\deg Z=0$.  Since $D$ is irreducible, its
intersections with distinct curves are nonnegative; together with
$D^2=0$, this proves nefness.

Write $E_j$ for the exceptional curves.  The blowup formulas give
\[
 K_S=\pi^*K_{S_0}+\sum_jE_j,\qquad
 D=\pi^*C-\sum_jE_j.
\]
Thus $K_S+D\sim\pi^*(K_{S_0}+C)$, proving the displayed formula.
For $(a,b)=(2,b)$ with $b\ge3$, this is the pullback of
$\OO_{\PP^1}(b-2)$ under the second projection.  For $a,b\ge3$
the corresponding line bundle on $S_0$ is ample.  In both cases its
pullback is nef and semiample, with the claimed Iitaka dimension.

If $N$ is non-torsion, then $H^0(D,N^{\otimes m})=0$ for every
$m>0$.  Induction in the sequences
\[
 0\longrightarrow\OO_S((m-1)D)
 \longrightarrow\OO_S(mD)
 \longrightarrow N^{\otimes m}\longrightarrow0
\]
gives $h^0(S,\OO_S(mD))=1$ for all $m\ge0$.  This is exactly
$\kappa(S,D)=0$.
\end{proof}

\begin{remark}[Crepant boundary-point blowups]
\label{rem:surface-crepant-blowup}
The exceptional curves in
Proposition~\ref{prop:surface-arbitrary-normal} satisfy
$(K_S+D)\cdot E_j=0$.  Thus they are not removed by a negative
log extremal contraction.  Moreover $E_j\setminus D\simeq\C$
belongs to the open complement.  Blowing it down sends this open
curve into a point of the removed boundary and is not a proper
modification of $S\setminus D$ onto $S_0\setminus C$.  This
distinguishes these blowups from boundary-node blowups with full
reduced total transform, which preserve the open complement.
\end{remark}

\section{Holonomy obstructions and the remaining surface problem}
\label{sec:surface-remaining}

The previous sections give complete criteria in the nonprojective and
smooth ruled settings, as well as further affirmative projective
cases. We now isolate what the neighborhood arguments imply, and what
is still missing for an arbitrary continuously weakly complete surface.

\subsection{Two obstructions for foliated neighborhoods}

The statements below concern the indicated foliated or suspension neighborhoods.
They neither assert that an arbitrary plurisubharmonic collar is basic nor
produce a foliation from the existence of such a collar.

\begin{proposition}[A continuous basic collar linearizes holonomy]
\label{prop:basic-collar}
Let $C$ be a compact leaf of a nonsingular holomorphic foliation on a complex
surface, defined on a neighborhood $U$ of $C$. Suppose that a finite continuous
function $u:U\setminus C\to\R$ is constant on leaves and tends to
$+\infty$ as the point approaches $C$. Then the holonomy along $C$ is
simultaneously holomorphically conjugate to a group of rotations.
On a possibly smaller neighborhood $V$, the divisor bundle $\OO_V(C)$ is
unitary flat. In particular, $V\setminus C$ admits a smooth pluriharmonic
function tending to $+\infty$ along $C$.
\end{proposition}

\begin{proof}
Fix a transverse coordinate $w$ with $C=\{w=0\}$, and put $v=u|_{\{w\ne0\}}$.
Take representatives of finitely many holonomy generators and their inverses.
Choose $0<r<R$ so that these maps are univalent on a neighborhood of
$\overline{B_R}$, carry $B_r$ into $B_R$, and preserve $v$ wherever used.
The closed disk $\overline{B_R}$ is chosen inside the transverse chart.
Take
\[
 t>\max_{r\le |w|\le R}v(w),\qquad
 E=\{0\}\cup\{w\in B_R\setminus\{0\}:v(w)>t\}.
\]
The divergence of $v$ makes $E$ open at $0$, and $E\Subset B_r$.
Invariance of $v$ gives $h(E)\subset E$ and $h^{-1}(E)\subset E$, hence
$h(E)=E$, for every chosen generator $h$. The component $D$ of $E$
containing $0$ is consequently a bounded domain preserved by all these maps.

The isotropy group $G=\operatorname{Aut}(D,0)$ is compact in the compact-open
topology. This standard normal-family fact can be seen by applying Cauchy's
estimate to all iterates and their inverses, which gives $|g'(0)|=1$, and then
taking simultaneous locally uniform limits of automorphisms and their inverses.
Let $\mu$ be normalized Haar measure on $G$ and define
\[
 \psi(w)=\int_G g'(0)^{-1}g(w)\,d\mu(g).
\]
This is holomorphic on $D$, satisfies $\psi(0)=0$ and $\psi'(0)=1$, and
right invariance of Haar measure gives
\[
 \psi(h(w))=h'(0)\psi(w)\qquad(h\in G).
\]
Thus $\psi$ is a local simultaneous linearizing coordinate and the multipliers
are unitary. No filling of complementary components of $D$ is required.

Transport this coordinate along finitely many flow boxes covering $C$.
Path changes give holonomy germs, so, after shrinking, the resulting
holomorphic first integrals $w_i$ satisfy
$w_i=\lambda_{ij}w_j$ on connected overlaps, with constant
$|\lambda_{ij}|=1$. They are local defining functions for $C$.
Their ratios therefore give unitary-constant transition functions for
$\OO_V(C)$ and define a smooth flat Hermitian metric on that bundle.
The functions $-\log|w_i|$ agree, are smooth and pluriharmonic off $C$,
and diverge along $C$.
\end{proof}

\begin{proposition}[Cyclic nonlinearizable suspensions have no psh collar]
\label{prop:cyclic-suspension}
Let $C$ be a compact connected Riemann surface of genus at least one, let
$\chi:\pi_1(C)\to\mathbb Z$ be primitive, and let $f$ be a nonlinearizable
irrationally indifferent holomorphic diffeomorphism germ at $0$.
Suppose that the germ of a neighborhood of $C$ is the holomorphic suspension
of $\gamma\mapsto f^{\chi(\gamma)}$.
There is no finite continuous plurisubharmonic function on a punctured
neighborhood of $C$ tending to $+\infty$ along $C$.
\end{proposition}

\begin{proof}
Choose an oriented nonseparating smooth simple closed curve $a$ whose
Poincar\'e dual is $\chi$. Cutting along $a$ gives a connected compact
bordered surface $B$ with two boundary circles. The character $\chi$ vanishes
on $\pi_1(B)$, and the cyclic cover $C_\chi$ is obtained by chaining copies
$B_n$, $n\in\mathbb Z$. The suspension is trivial over $B$; its two sides
are glued using $f$.

We use the following precise dynamical input: in arbitrarily small admissible
transverse neighborhoods there is a nontrivial compact connected completely
invariant hedgehog $K$, with $0\in\partial K$, and a point
$x\in\partial K$ whose forward orbit is dense in $\partial K$.
The existence and two-sided density are recalled in
\cite[Theorems~2.1 and~2.3]{KoikeOgawa};
the passage to forward density uses P\'erez-Marco's recurrence theorem,
as explicitly explained in \cite[Section~5]{KoikeOgawa}. Forward density is not being
inferred from two-sided density alone.

Choose $K$ small enough that its compact suspension is contained in the
neighborhood on which the proposed function $u$ is defined off $C$.
The point $x$ is nonperiodic, since its orbit closure contains both $0$ and
other points. The leaf through $x$ is therefore biholomorphic to $C_\chi$.
Write $w_n=f^n(x)$. Choose the boundary-coordinate convention so that the
cut preceding $B_n$ has parameter $w_n$, and the interior of $B_n$ is
the copy of $B$ with transverse parameter $w_n$.

Fix $y\in\partial K\setminus\{0\}$. The lifted cut circle with parameter
$y$ is compact and disjoint from $C$. Continuity gives a transverse
neighborhood $V_y$ of $y$ and a real number $M$ such that $u<M$ on all cut
circles with parameter in $V_y$. Fix a point $b$ in the interior of $B$.
Divergence along $C$ gives a neighborhood $V_0$ of $0$ for which
$u(b,w)>M$ whenever $w\in V_0\setminus\{0\}$.
Every tail of the forward dense orbit is dense: $\partial K$ has no isolated
points, so deleting finitely many orbit points cannot destroy density.
We may consequently choose
\[
 n_1<n_2<n_3,\qquad
 w_{n_1},w_{n_3}\in V_y,\qquad w_{n_2}\in V_0.
\]
The chain $B_{n_1}\cup\cdots\cup B_{n_3-1}$ is a compact connected
bordered Riemann surface in this leaf. Its only boundary components are the
cuts with parameters $w_{n_1}$ and $w_{n_3}$, where $u<M$; the interior
point $(b,w_{n_2})$ has $u>M$. The restriction of $u$ to the leaf is
subharmonic, contradicting the maximum principle.
\end{proof}

\begin{remark}[Scope]
Subharmonicity is unnecessary in Proposition~\ref{prop:basic-collar}, but
leafwise constancy is essential to its argument. Proposition~\ref{prop:cyclic-suspension}
does not require leafwise constancy or a parabolic universal cover; it does
require the specified holomorphic cyclic suspension. Neither proposition
upgrades an arbitrary continuous psh collar to a basic one or supplies a
foliation or holomorphic suspension for a general surface neighborhood.
Consequently these obstructions do not settle the unrestricted
continuous-to-smooth implication.
\end{remark}


\begin{remark}[Pseudoconvex neighborhoods are a different condition]
\label{rem:pseudoconvex-neighborhoods}
Koike--Takakura \cite[Theorem~1.1]{KoikeTakakura2026} consider a
holomorphic section of a holomorphic submersion with one-dimensional
fibers over a compact complex manifold, with unitary flat normal
bundle. They prove that a fundamental system of locally pseudoconvex
neighborhoods is equivalent to a holomorphic tubular neighborhood.
The submersion and the pseudoconvex neighborhood basis are hypotheses
of that result, not consequences of the collar assumption here.
In particular, a smooth psh function $u$ tending to $+\infty$ along
$C$ gives, at regular levels, neighborhoods on the side $\{u>t\}$
with local defining function $t-u$. Its Levi form on complex tangent
directions is nonpositive. Thus a divergent psh function on the
\emph{punctured} neighborhood does not by itself provide the
pseudoconvex neighborhood basis required in that theorem. We do not
use it to close the remaining surface argument.
\end{remark}

\subsection{Necessary conditions in the smooth-boundary case}

\begin{corollary}[Necessary conditions for a smooth-boundary counterexample]
\label{cor:surface-residual-homology}
Suppose that $S$ is a smooth projective surface, $C\subset S$ is a smooth connected curve
with $C^2=0$, and $X=S\setminus C$ has a continuous plurisubharmonic exhaustion but no
smooth one. Then
\[
 \OO(X)=\C,\qquad q(S)<g(C),
\]
and
\[
 h^0(S,\OO_S(mC))=1\qquad(m\ge1),
\]
and the unitary character of $N_{C/S}$ has infinite image on
$K=\ker(H_1(C,\mathbb Z)\to H_1(S,\mathbb Z))$. Moreover,
\[
 \operatorname{rank}K=2\bigl(g(C)-q(S)\bigr)\ge2.
\]
The Ueda type is infinite, the neighborhood is not unitary flat, and $N_{C/S}$ is
non-torsion and fails the Diophantine condition in
Proposition~\ref{prop:surface-ueda-sufficient}. Using the September 2026 preprint
\cite[Theorem 1.6]{CFSTZ2026}, one further has $g(C)\ge2$.
In addition, every continuous psh exhaustion $u$ of $X$ satisfies
$(a_u,+\infty)\subset u(\Sigma_X^0)$ for some $a_u\in\R$.
\end{corollary}
\begin{proof}
Proposition~\ref{prop:surface-nonconstant-function} gives $\OO(X)=\C$. For every section $s$
of $\OO_S(mC)$, the quotient $s/s_C^m$ is holomorphic on $X$, hence constant. The canonical
section $s_C^m$ is nonzero, so $h^0(S,\OO_S(mC))=1$. Corollary~\ref{cor:irregularity} and
Theorem~\ref{thm:homology} give the homological conditions. If the restriction
$H^1(S,\OO_S)\to H^1(C,\OO_C)$ were not injective, Proposition~2.7 of \cite{CLPT2018} would
make $C$ a multiple fiber, again giving a smooth exhaustion. It is therefore injective.
Hodge decomposition and duality make $H_1(C,\mathbb Q)\to H_1(S,\mathbb Q)$ surjective.
Their dimensions are $2g(C)$ and $2q(S)$, respectively. Tensoring the integral kernel
sequence with $\mathbb Q$ proves the stated rank formula, even when the target has torsion.
The remaining exclusions follow from Proposition~\ref{prop:surface-ueda-sufficient}. The
minimal-kernel assertion is
Proposition~\ref{prop:surface-minimal-kernel}. The elliptic exclusion
uses the stated 2026 preprint; the homological and irregularity criteria
do not depend on that preprint.
\end{proof}


\subsection{The remaining analytic problem}

\begin{remark}[The remaining smooth-boundary subproblem]
\label{rem:surface-unproved-collar}
Even when $C$ is a smooth connected curve in a smooth projective
surface $S$, with $C^2=0$ and $g(C)\ge2$, the following implication
has not been proved here:
\[
 \begin{gathered}
 u\in C^0(U\setminus C)\cap\operatorname{PSH}(U\setminus C),
 \qquad \lim_{x\to C}u(x)=+\infty
 \\[2pt]
 \Longrightarrow\quad
 \exists\,\Phi\in C^\infty(V\setminus C)\cap\operatorname{PSH}(V\setminus C),
 \qquad \lim_{x\to C}\Phi(x)=+\infty,
 \end{gathered}
\]
where $U$ and $V$ are neighborhoods of $C$.  A smooth collar would
extend to a smooth exhaustion by Lemma~\ref{lem:surface-collar-gluing}.
Proposition~\ref{prop:surface-ueda-sufficient} settles finite Ueda
type, unitary flat neighborhoods, torsion or Diophantine normal bundles,
the stated monodromy-extension condition, and elliptic boundary curves.
The homological and irregularity criteria sharpen this class as
in Corollary~\ref{cor:surface-residual-homology}. The remaining
smooth-boundary situation has infinite Ueda type and non-torsion
normal monodromy with infinite image on the homology kernel. Knowing
only the degree-zero normal bundle does not settle it. The foliated results above do not
remove the gap: an arbitrary
continuous psh collar need not have been shown to be basic, and a
holomorphic suspension is additional geometric structure. This article
does not prove the unrestricted continuous-to-smooth implication for surfaces.

Reducible weighted square-zero boundaries and singular log endpoints
with changed boundary coefficients require additional arguments as
well. Semiampleness of $K_S+D$ provides a map of the compact pair, but
does not ensure that its restriction to the divisor complement is proper.

In particular, an almost-positive regularization estimate is
insufficient on its own.  If
$i\partial\bar\partial u_\varepsilon\ge-\varepsilon\omega$, then
\[
 i\partial\bar\partial\chi(u_\varepsilon)
 =\chi'(u_\varepsilon)i\partial\bar\partial u_\varepsilon
 +\chi''(u_\varepsilon)i\partial u_\varepsilon
       \wedge\bar\partial u_\varepsilon.
\]
The second term controls only one complex direction.  On
$\ker\partial u_\varepsilon$, a negative Levi direction can remain.
Therefore convex composition alone does not turn such estimates into
the missing semipositive exhaustion.
\end{remark}


\subsection{Conclusion and dependence on recent results}

The threefold counterexamples distinguish continuous weak completeness
from $C^2$ weak completeness, including in the quasi-projective
K\"ahler category and after a modification supported over one point.
In dimension two, the numerical criteria completely settle divisor
complements with nonprojective compactifier and reduced SNC complements
in a smooth ruling. The projective residual problem still includes
connected weighted square-zero boundaries and the smooth-boundary
class in Corollary~\ref{cor:surface-residual-homology}.

The elliptic and unitary-period rational-cycle arguments explicitly
use \cite{CFSTZ2026}, a preprint submitted September~22, 2026.
Section~1.3.1 of that preprint records the construction of a
type~$(\gamma)$ embedding in a compact complex surface as an open
problem, and Theorem~1.6 excludes elliptic curves. A smooth-boundary
counterexample in our residual class would give such a compact
type~$(\gamma)$ embedding. The converse is not established: such an
embedding alone would give neither a continuous psh exhaustion of its
complement nor a proof that every smooth psh exhaustion is impossible.
Neither the homological criterion nor the all-degree smooth-ruling
criterion depends on this preprint.

Finally, weak completeness has not supplied a divisorial
compactification of an arbitrary noncompact surface. Accordingly the
full dimension-two question remains unresolved by the arguments here,
both without further assumptions and under the sole additional
assumption of K\"ahlerness or quasi-projectivity.

\begin{thebibliography}{99}
\raggedright

\bibitem{Koike2021}
T.~Koike,
\emph{Linearization of transition functions of a semi-positive line
bundle along a certain submanifold},
Ann. Inst. Fourier \textbf{71} (2021), no.~5, 2237--2271.
\href{https://doi.org/10.5802/aif.3439}{doi:10.5802/aif.3439}.

\bibitem{Wu2026}
X.~Wu,
\emph{Generalized Ueda Obstruction Classes and Applications to
Non-semi-positivity of Line Bundles},
preprint (2026), arXiv:2606.22821v1.
\href{https://arxiv.org/abs/2606.22821v1}{arXiv:2606.22821}.

\bibitem{Zhang2025}
Y.~Zhang,
\emph{On the Non-Semipositivity of a Nef and Big Line Bundle on
Grauert's Example},
preprint (2025), arXiv:2512.23446v1.
\href{https://arxiv.org/abs/2512.23446v1}{arXiv:2512.23446}.

\bibitem{Ohsawa1981}
T.~Ohsawa,
\emph{Weakly $1$-complete manifold and Levi problem},
Publ. Res. Inst. Math. Sci. \textbf{17} (1981), no.~1, 153--164.
\href{https://doi.org/10.2977/PRIMS/1195186709}{doi:10.2977/PRIMS/1195186709}.
\url{https://ems.press/content/serial-article-files/42043}.

\bibitem{TomassiniVajaitu2008}
G.~Tomassini and V.~V\^aj\^aitu,
\emph{Weakly $1$-complete surfaces with singularities and applications},
Michigan Math. J. \textbf{56} (2008), no.~3, 483--494.
\href{https://doi.org/10.1307/mmj/1231770355}{doi:10.1307/mmj/1231770355}.

\bibitem{Vajaitu2022}
V.~V\^aj\^aitu,
\emph{Holomorphic convexity and pseudoconvex surfaces},
Ann. Sc. Norm. Super. Pisa Cl. Sci. (5) \textbf{23} (2022),
no.~4, 1659--1678.
\href{https://doi.org/10.2422/2036-2145.202011_018}{doi:10.2422/2036-2145.202011\_018}.
\url{https://journals.sns.it/index.php/annaliscienze/article/view/5700}.

\bibitem{MongodiSlodkowski2020}
S.~Mongodi and Z.~Slodkowski,
\emph{Domains with a continuous exhaustion in weakly complete surfaces},
Math. Z. \textbf{296} (2020), nos.~3--4, 1011--1019.
\href{https://doi.org/10.1007/s00209-020-02466-z}{doi:10.1007/s00209-020-02466-z}.
\url{https://arxiv.org/abs/1904.03497}.

\bibitem{Grauert1962}
H.~Grauert,
\emph{\"Uber Modifikationen und exzeptionelle analytische Mengen},
Math. Ann. \textbf{146} (1962), 331--368.
\href{https://doi.org/10.1007/BF01441136}{doi:10.1007/BF01441136}.

\bibitem{FN2005}
Y.~Fujimoto and N.~Nakayama,
\emph{Compact complex surfaces admitting non-trivial surjective endomorphisms},
Tohoku Math. J. (2) \textbf{57} (2005), no.~3, 395--426.
\href{https://doi.org/10.2748/tmj/1128703004}{doi:10.2748/tmj/1128703004}.
\href{https://www.kurims.kyoto-u.ac.jp/preprint/file/RIMS1431.pdf}{Author preprint, Section~2 and Lemma~2.1}.

\bibitem{Enoki}
I.~Enoki, \emph{Surfaces of class $\mathrm{VII}_0$ with curves}, Tohoku Math. J. (2)
\textbf{33} (1981), no.~4, 453--492.
\href{https://doi.org/10.2748/tmj/1178229349}{doi:10.2748/tmj/1178229349}.
\href{https://www.jstage.jst.go.jp/article/tmj1949/33/4/33_4_453/_pdf/-char/en}{Original article, pp.453--454}.

\bibitem{Fujino2012}
O.~Fujino,
\emph{Minimal model theory for log surfaces},
Publ. Res. Inst. Math. Sci. \textbf{48} (2012), no.~2, 339--371.
\href{https://doi.org/10.2977/PRIMS/71}{doi:10.2977/PRIMS/71}.

\bibitem{Fujino2021}
O.~Fujino,
\emph{Minimal model theory for log surfaces in Fujiki's class $\mathcal C$},
Nagoya Math. J. \textbf{244} (2021), 256--282.
\href{https://doi.org/10.1017/nmj.2020.14}{doi:10.1017/nmj.2020.14}.

\bibitem{Ueda1983}
T.~Ueda,
\emph{On the neighborhood of a compact complex curve with topologically
trivial normal bundle},
J. Math. Kyoto Univ. \textbf{22} (1982/83), no.~4, 583--607.
\href{https://doi.org/10.1215/kjm/1250521670}{doi:10.1215/kjm/1250521670}.

\bibitem{CLPT2018}
B.~Claudon, F.~Loray, J.~V.~Pereira and F.~Touzet,
\emph{Compact leaves of codimension one holomorphic foliations on
projective manifolds},
Ann. Sci. \'Ecole Norm. Sup. (4) \textbf{51} (2018), no.~6, 1457--1506.
\href{https://doi.org/10.24033/asens.2379}{doi:10.24033/asens.2379}.
\href{https://arxiv.org/abs/1512.06623}{Author preprint, Theorems~1.3, 2.5, 4.4 and Proposition~2.7}.

\bibitem{KoikeNodes2017}
T.~Koike,
\emph{Ueda theory for compact curves with nodes},
Indiana Univ. Math. J. \textbf{66} (2017), no.~3, 845--876.
\href{https://arxiv.org/abs/1507.00109}{arXiv:1507.00109}.

\bibitem{CFSTZ2026}
Y.~Chen, S.~Filip, S.~Sun, V.~Tosatti and J.~Zhang,
\emph{Compact K\"ahler surfaces with semipositive anticanonical bundle},
preprint (2026), arXiv:2609.26716v1, submitted September~22, 2026.
\href{https://arxiv.org/abs/2609.26716v1}{arXiv:2609.26716}.

\bibitem{KoikeUeda2025}
T.~Koike and T.~Ueda,
\emph{Plurisubharmonic functions on affine line bundles over compact
K\"ahler manifolds},
Publ. Res. Inst. Math. Sci. \textbf{61} (2025), no.~1, 139--152.
\href{https://doi.org/10.4171/PRIMS/61-1-3}{doi:10.4171/PRIMS/61-1-3}.



\bibitem{KoikeOgawa}
T.~Koike and N.~Ogawa,
\emph{On the neighborhood of a torus leaf and dynamics of holomorphic foliations},
Tohoku Math. J. (2) \textbf{77} (2025), no.~3, 421--443.
\href{https://doi.org/10.2748/tmj.20231108}{doi:10.2748/tmj.20231108}.
\href{https://arxiv.org/abs/1808.10219v3}{Author preprint, arXiv:1808.10219v3}.
The theorem numbering used for the dynamical inputs refers to
Theorems~2.1, 2.3 and Section~5 of this author-preprint version.

\bibitem{KoikeTakakura2026}
T.~Koike and M.~Takakura,
\emph{Rigidity of Pseudoconvex Neighborhoods and Ohsawa--Takegoshi Extension},
preprint (2026), arXiv:2610.29082v1, submitted October~21, 2026.
\href{https://arxiv.org/abs/2610.29082v1}{arXiv:2610.29082v1}.

\bibitem{Barlet2015}
D.~Barlet,
\emph{Some examples due to H.~Hironaka},
preprint (2015), arXiv:1501.01129v2, Section~6.
\href{https://arxiv.org/abs/1501.01129}{arXiv:1501.01129}.

\bibitem{StacksHironaka}
The Stacks Project Authors,
\emph{The Stacks Project},
Lemma~110.67.1, Tag~08J1, ``Being projective is not local on the base.''
\url{https://stacks.math.columbia.edu/tag/08J1}.

\end{thebibliography}
\end{document}
